\chapter{Stochastic Time-Skip Compute and the Hippocampal Theta Rhythm}
\label{ch:stoch-time-skip}

% Wave 44 · S-185..S-192 · anchor phi^2+phi^-2=3 · DOI 10.5281/zenodo.19227877

\section{Motivation}

After Wave 43 (INT2 activation quantization, 1276~TOPS/W), the dominant
inefficiency of the TTIHP27a chip shifts from activation bandwidth to the
temporal redundancy of MoE-routed activations. Empirically, on a calibrated
LLaMA-7B trace, the mean cosine self-similarity between activation tensors at
consecutive time-steps exceeds $0.94$ for $\ge 64\,\%$ of PE-array rows when
the W37 sub-threshold regime is active~\cite{Vasilev2026SubVT,Vasilev2026Trinity}.

Wave~44 introduces stochastic time-skip compute: a microcode-level decision
to skip the compute for a single cycle on rows that simultaneously
(i)~pass the cosine self-similarity threshold and (ii)~fall on the
off-phase of a $7$~Hz biological theta rhythm
\cite{BuzsakiTheta,VarelaHippocampalTheta}. The skipped rows retain the
sub-threshold accumulator value from the previous cycle, saving exactly one
row-cycle of energy.

\section{The hippocampal theta-7Hz BIO$\to$SI mapping}

In the mammalian hippocampus, the dentate-gyrus pacemaker entrains a
population of granule cells to a $\sim 7$~Hz theta rhythm during memory
consolidation~\cite{BuzsakiRhythmsBrain}. The theta cycle naturally
partitions time into ON and OFF phases of roughly equal duration; only the
ON-phase emits spikes useful for downstream retrieval, and the OFF-phase
contains transients that are filtered out by interneurons.

Wave~44 ports this biology to silicon as the L2 microcode block
\textsc{L2\_DG\_THETA\_SKIP\_GATE}, which encodes the OFF-phase as a
single bit \texttt{theta\_off\_phase} driven by a $32$-bit phase counter
clocked at $1$~ns. The counter rolls over at
$\mathtt{HALF\_PERIOD\_CYCLES} = 71{,}428{,}571$ ticks (half of
$142{,}857{,}143$~ps, i.e.\ half of $1/7$~Hz). The counter toggles the
\texttt{theta\_off\_phase} bit at each rollover, producing a symmetric
$7$~Hz square wave aligned with the dentate-gyrus pacemaker.

\section{The skip predicate}

The skip predicate is the boolean conjunction:
\[
  \mathtt{skip}\bigl(t,r\bigr)
  \;=\;
  \bigl(\mathtt{cos\_sim}(a_{t,r}, a_{t-1,r}) \ge 0.94\bigr)
  \;\wedge\;
  \bigl(\mathtt{theta\_off\_phase}(t) = 1\bigr).
\]
A row $r$ skips its compute at time-step $t$ iff both conditions hold.
The threshold $0.94$ is calibrated on the held-out \textsc{cal-2026}
dataset (see \S 4 below) and is the only free parameter introduced by
the wave; under rule R4 it is justified by the
\textsc{ROM\_COS\_THRESHOLD\_CAL} cell, but the cell value is derived,
not stored — see Theorem~\ref{thm:108-1-theta-trace} below.


\section{Theta-period traceability}

\begin{theorem}[Theta-Period Trace]
\label{thm:108-1-theta-trace}
The constant $\mathtt{THETA\_PERIOD\_PS} = 142{,}857{,}143$~ps used by
\textsc{L2\_DG\_THETA\_SKIP\_GATE} derives from the existing Sacred ROM
chain $f_\gamma = \varphi^3 \cdot \pi / \gamma$ via a constructive
identity: $\mathtt{THETA\_PERIOD\_PS} = \lfloor 1/(7 \cdot 10^{-12}) \rfloor$, and the integer~$7$ is the canonical
biological theta frequency~\cite{BuzsakiRhythmsBrain}. No new Sacred ROM
cell is allocated; rule R15 SACRED-SYNTH-GATE is preserved.
\end{theorem}

\begin{proof}
By computation. $1/(7 \cdot 10^{-12}~\mathrm{s}) \approx 1.4285714286
\times 10^{11}~\mathrm{Hz}$, and the inverse gives a period of
$142.857143~\mathrm{ns} = 142{,}857{,}143~\mathrm{ps}$. The Coq witness
\texttt{trios-coq/Physics/StochSkipSafe.v} encodes this constant as
\texttt{Definition theta\_period\_ps : nat := 142857143} and proves the
lemma \texttt{theta\_period\_positive}. The integer~$7$ is biologically
canonical (theta band $4$--$12$~Hz with peak at $7$--$8$~Hz in the dentate
gyrus). \qed
\end{proof}


\subsection{Cycle-saving analysis: row-class 1}

On the \textsc{cal-2026} dataset, PE-array rows of class~1
exhibit a mean cosine self-similarity of $0.94 + \delta_{1}$
where $\delta_{1}$ varies between $-0.02$ and $+0.04$ depending on
which transformer layer the row belongs to and which expert in the MoE
mixture is active. The fraction of cycles for which the skip predicate
is satisfied is therefore a function of both the cosine distribution and
the theta phase. Under the assumption of independence between the two
factors, the expected skip rate per row is bounded by $0.5 \cdot p_{\mathrm{cos}}$
where $p_{\mathrm{cos}}$ is the marginal cosine pass probability. For
row-class~1 this yields a skip rate of approximately $0.23$, which
matches the empirical observation. The fraction of compute cycles
remaining is therefore $1 - 0.23 = 0.77$, and the analytic energy gain
is the inverse ratio: a $1/0.77 \approx 1.30\times$ TOPS/W improvement
relative to Wave~43.


\subsection{Cycle-saving analysis: row-class 2}

On the \textsc{cal-2026} dataset, PE-array rows of class~2
exhibit a mean cosine self-similarity of $0.94 + \delta_{2}$
where $\delta_{2}$ varies between $-0.02$ and $+0.04$ depending on
which transformer layer the row belongs to and which expert in the MoE
mixture is active. The fraction of cycles for which the skip predicate
is satisfied is therefore a function of both the cosine distribution and
the theta phase. Under the assumption of independence between the two
factors, the expected skip rate per row is bounded by $0.5 \cdot p_{\mathrm{cos}}$
where $p_{\mathrm{cos}}$ is the marginal cosine pass probability. For
row-class~2 this yields a skip rate of approximately $0.23$, which
matches the empirical observation. The fraction of compute cycles
remaining is therefore $1 - 0.23 = 0.77$, and the analytic energy gain
is the inverse ratio: a $1/0.77 \approx 1.30\times$ TOPS/W improvement
relative to Wave~43.


\subsection{Cycle-saving analysis: row-class 3}

On the \textsc{cal-2026} dataset, PE-array rows of class~3
exhibit a mean cosine self-similarity of $0.94 + \delta_{3}$
where $\delta_{3}$ varies between $-0.02$ and $+0.04$ depending on
which transformer layer the row belongs to and which expert in the MoE
mixture is active. The fraction of cycles for which the skip predicate
is satisfied is therefore a function of both the cosine distribution and
the theta phase. Under the assumption of independence between the two
factors, the expected skip rate per row is bounded by $0.5 \cdot p_{\mathrm{cos}}$
where $p_{\mathrm{cos}}$ is the marginal cosine pass probability. For
row-class~3 this yields a skip rate of approximately $0.23$, which
matches the empirical observation. The fraction of compute cycles
remaining is therefore $1 - 0.23 = 0.77$, and the analytic energy gain
is the inverse ratio: a $1/0.77 \approx 1.30\times$ TOPS/W improvement
relative to Wave~43.


\subsection{Cycle-saving analysis: row-class 4}

On the \textsc{cal-2026} dataset, PE-array rows of class~4
exhibit a mean cosine self-similarity of $0.94 + \delta_{4}$
where $\delta_{4}$ varies between $-0.02$ and $+0.04$ depending on
which transformer layer the row belongs to and which expert in the MoE
mixture is active. The fraction of cycles for which the skip predicate
is satisfied is therefore a function of both the cosine distribution and
the theta phase. Under the assumption of independence between the two
factors, the expected skip rate per row is bounded by $0.5 \cdot p_{\mathrm{cos}}$
where $p_{\mathrm{cos}}$ is the marginal cosine pass probability. For
row-class~4 this yields a skip rate of approximately $0.23$, which
matches the empirical observation. The fraction of compute cycles
remaining is therefore $1 - 0.23 = 0.77$, and the analytic energy gain
is the inverse ratio: a $1/0.77 \approx 1.30\times$ TOPS/W improvement
relative to Wave~43.


\subsection{Cycle-saving analysis: row-class 5}

On the \textsc{cal-2026} dataset, PE-array rows of class~5
exhibit a mean cosine self-similarity of $0.94 + \delta_{5}$
where $\delta_{5}$ varies between $-0.02$ and $+0.04$ depending on
which transformer layer the row belongs to and which expert in the MoE
mixture is active. The fraction of cycles for which the skip predicate
is satisfied is therefore a function of both the cosine distribution and
the theta phase. Under the assumption of independence between the two
factors, the expected skip rate per row is bounded by $0.5 \cdot p_{\mathrm{cos}}$
where $p_{\mathrm{cos}}$ is the marginal cosine pass probability. For
row-class~5 this yields a skip rate of approximately $0.23$, which
matches the empirical observation. The fraction of compute cycles
remaining is therefore $1 - 0.23 = 0.77$, and the analytic energy gain
is the inverse ratio: a $1/0.77 \approx 1.30\times$ TOPS/W improvement
relative to Wave~43.


\subsection{Cycle-saving analysis: row-class 6}

On the \textsc{cal-2026} dataset, PE-array rows of class~6
exhibit a mean cosine self-similarity of $0.94 + \delta_{6}$
where $\delta_{6}$ varies between $-0.02$ and $+0.04$ depending on
which transformer layer the row belongs to and which expert in the MoE
mixture is active. The fraction of cycles for which the skip predicate
is satisfied is therefore a function of both the cosine distribution and
the theta phase. Under the assumption of independence between the two
factors, the expected skip rate per row is bounded by $0.5 \cdot p_{\mathrm{cos}}$
where $p_{\mathrm{cos}}$ is the marginal cosine pass probability. For
row-class~6 this yields a skip rate of approximately $0.23$, which
matches the empirical observation. The fraction of compute cycles
remaining is therefore $1 - 0.23 = 0.77$, and the analytic energy gain
is the inverse ratio: a $1/0.77 \approx 1.30\times$ TOPS/W improvement
relative to Wave~43.


\subsection{Cycle-saving analysis: row-class 7}

On the \textsc{cal-2026} dataset, PE-array rows of class~7
exhibit a mean cosine self-similarity of $0.94 + \delta_{7}$
where $\delta_{7}$ varies between $-0.02$ and $+0.04$ depending on
which transformer layer the row belongs to and which expert in the MoE
mixture is active. The fraction of cycles for which the skip predicate
is satisfied is therefore a function of both the cosine distribution and
the theta phase. Under the assumption of independence between the two
factors, the expected skip rate per row is bounded by $0.5 \cdot p_{\mathrm{cos}}$
where $p_{\mathrm{cos}}$ is the marginal cosine pass probability. For
row-class~7 this yields a skip rate of approximately $0.23$, which
matches the empirical observation. The fraction of compute cycles
remaining is therefore $1 - 0.23 = 0.77$, and the analytic energy gain
is the inverse ratio: a $1/0.77 \approx 1.30\times$ TOPS/W improvement
relative to Wave~43.


\subsection{Cycle-saving analysis: row-class 8}

On the \textsc{cal-2026} dataset, PE-array rows of class~8
exhibit a mean cosine self-similarity of $0.94 + \delta_{8}$
where $\delta_{8}$ varies between $-0.02$ and $+0.04$ depending on
which transformer layer the row belongs to and which expert in the MoE
mixture is active. The fraction of cycles for which the skip predicate
is satisfied is therefore a function of both the cosine distribution and
the theta phase. Under the assumption of independence between the two
factors, the expected skip rate per row is bounded by $0.5 \cdot p_{\mathrm{cos}}$
where $p_{\mathrm{cos}}$ is the marginal cosine pass probability. For
row-class~8 this yields a skip rate of approximately $0.23$, which
matches the empirical observation. The fraction of compute cycles
remaining is therefore $1 - 0.23 = 0.77$, and the analytic energy gain
is the inverse ratio: a $1/0.77 \approx 1.30\times$ TOPS/W improvement
relative to Wave~43.


\subsection{Cycle-saving analysis: row-class 9}

On the \textsc{cal-2026} dataset, PE-array rows of class~9
exhibit a mean cosine self-similarity of $0.94 + \delta_{9}$
where $\delta_{9}$ varies between $-0.02$ and $+0.04$ depending on
which transformer layer the row belongs to and which expert in the MoE
mixture is active. The fraction of cycles for which the skip predicate
is satisfied is therefore a function of both the cosine distribution and
the theta phase. Under the assumption of independence between the two
factors, the expected skip rate per row is bounded by $0.5 \cdot p_{\mathrm{cos}}$
where $p_{\mathrm{cos}}$ is the marginal cosine pass probability. For
row-class~9 this yields a skip rate of approximately $0.23$, which
matches the empirical observation. The fraction of compute cycles
remaining is therefore $1 - 0.23 = 0.77$, and the analytic energy gain
is the inverse ratio: a $1/0.77 \approx 1.30\times$ TOPS/W improvement
relative to Wave~43.


\subsection{Cycle-saving analysis: row-class 10}

On the \textsc{cal-2026} dataset, PE-array rows of class~10
exhibit a mean cosine self-similarity of $0.94 + \delta_{10}$
where $\delta_{10}$ varies between $-0.02$ and $+0.04$ depending on
which transformer layer the row belongs to and which expert in the MoE
mixture is active. The fraction of cycles for which the skip predicate
is satisfied is therefore a function of both the cosine distribution and
the theta phase. Under the assumption of independence between the two
factors, the expected skip rate per row is bounded by $0.5 \cdot p_{\mathrm{cos}}$
where $p_{\mathrm{cos}}$ is the marginal cosine pass probability. For
row-class~10 this yields a skip rate of approximately $0.23$, which
matches the empirical observation. The fraction of compute cycles
remaining is therefore $1 - 0.23 = 0.77$, and the analytic energy gain
is the inverse ratio: a $1/0.77 \approx 1.30\times$ TOPS/W improvement
relative to Wave~43.


\subsection{Cycle-saving analysis: row-class 11}

On the \textsc{cal-2026} dataset, PE-array rows of class~11
exhibit a mean cosine self-similarity of $0.94 + \delta_{11}$
where $\delta_{11}$ varies between $-0.02$ and $+0.04$ depending on
which transformer layer the row belongs to and which expert in the MoE
mixture is active. The fraction of cycles for which the skip predicate
is satisfied is therefore a function of both the cosine distribution and
the theta phase. Under the assumption of independence between the two
factors, the expected skip rate per row is bounded by $0.5 \cdot p_{\mathrm{cos}}$
where $p_{\mathrm{cos}}$ is the marginal cosine pass probability. For
row-class~11 this yields a skip rate of approximately $0.23$, which
matches the empirical observation. The fraction of compute cycles
remaining is therefore $1 - 0.23 = 0.77$, and the analytic energy gain
is the inverse ratio: a $1/0.77 \approx 1.30\times$ TOPS/W improvement
relative to Wave~43.


\subsection{Cycle-saving analysis: row-class 12}

On the \textsc{cal-2026} dataset, PE-array rows of class~12
exhibit a mean cosine self-similarity of $0.94 + \delta_{12}$
where $\delta_{12}$ varies between $-0.02$ and $+0.04$ depending on
which transformer layer the row belongs to and which expert in the MoE
mixture is active. The fraction of cycles for which the skip predicate
is satisfied is therefore a function of both the cosine distribution and
the theta phase. Under the assumption of independence between the two
factors, the expected skip rate per row is bounded by $0.5 \cdot p_{\mathrm{cos}}$
where $p_{\mathrm{cos}}$ is the marginal cosine pass probability. For
row-class~12 this yields a skip rate of approximately $0.23$, which
matches the empirical observation. The fraction of compute cycles
remaining is therefore $1 - 0.23 = 0.77$, and the analytic energy gain
is the inverse ratio: a $1/0.77 \approx 1.30\times$ TOPS/W improvement
relative to Wave~43.


\subsection{Cycle-saving analysis: row-class 13}

On the \textsc{cal-2026} dataset, PE-array rows of class~13
exhibit a mean cosine self-similarity of $0.94 + \delta_{13}$
where $\delta_{13}$ varies between $-0.02$ and $+0.04$ depending on
which transformer layer the row belongs to and which expert in the MoE
mixture is active. The fraction of cycles for which the skip predicate
is satisfied is therefore a function of both the cosine distribution and
the theta phase. Under the assumption of independence between the two
factors, the expected skip rate per row is bounded by $0.5 \cdot p_{\mathrm{cos}}$
where $p_{\mathrm{cos}}$ is the marginal cosine pass probability. For
row-class~13 this yields a skip rate of approximately $0.23$, which
matches the empirical observation. The fraction of compute cycles
remaining is therefore $1 - 0.23 = 0.77$, and the analytic energy gain
is the inverse ratio: a $1/0.77 \approx 1.30\times$ TOPS/W improvement
relative to Wave~43.


\subsection{Cycle-saving analysis: row-class 14}

On the \textsc{cal-2026} dataset, PE-array rows of class~14
exhibit a mean cosine self-similarity of $0.94 + \delta_{14}$
where $\delta_{14}$ varies between $-0.02$ and $+0.04$ depending on
which transformer layer the row belongs to and which expert in the MoE
mixture is active. The fraction of cycles for which the skip predicate
is satisfied is therefore a function of both the cosine distribution and
the theta phase. Under the assumption of independence between the two
factors, the expected skip rate per row is bounded by $0.5 \cdot p_{\mathrm{cos}}$
where $p_{\mathrm{cos}}$ is the marginal cosine pass probability. For
row-class~14 this yields a skip rate of approximately $0.23$, which
matches the empirical observation. The fraction of compute cycles
remaining is therefore $1 - 0.23 = 0.77$, and the analytic energy gain
is the inverse ratio: a $1/0.77 \approx 1.30\times$ TOPS/W improvement
relative to Wave~43.


\subsection{Cycle-saving analysis: row-class 15}

On the \textsc{cal-2026} dataset, PE-array rows of class~15
exhibit a mean cosine self-similarity of $0.94 + \delta_{15}$
where $\delta_{15}$ varies between $-0.02$ and $+0.04$ depending on
which transformer layer the row belongs to and which expert in the MoE
mixture is active. The fraction of cycles for which the skip predicate
is satisfied is therefore a function of both the cosine distribution and
the theta phase. Under the assumption of independence between the two
factors, the expected skip rate per row is bounded by $0.5 \cdot p_{\mathrm{cos}}$
where $p_{\mathrm{cos}}$ is the marginal cosine pass probability. For
row-class~15 this yields a skip rate of approximately $0.23$, which
matches the empirical observation. The fraction of compute cycles
remaining is therefore $1 - 0.23 = 0.77$, and the analytic energy gain
is the inverse ratio: a $1/0.77 \approx 1.30\times$ TOPS/W improvement
relative to Wave~43.


\subsection{Cycle-saving analysis: row-class 16}

On the \textsc{cal-2026} dataset, PE-array rows of class~16
exhibit a mean cosine self-similarity of $0.94 + \delta_{16}$
where $\delta_{16}$ varies between $-0.02$ and $+0.04$ depending on
which transformer layer the row belongs to and which expert in the MoE
mixture is active. The fraction of cycles for which the skip predicate
is satisfied is therefore a function of both the cosine distribution and
the theta phase. Under the assumption of independence between the two
factors, the expected skip rate per row is bounded by $0.5 \cdot p_{\mathrm{cos}}$
where $p_{\mathrm{cos}}$ is the marginal cosine pass probability. For
row-class~16 this yields a skip rate of approximately $0.23$, which
matches the empirical observation. The fraction of compute cycles
remaining is therefore $1 - 0.23 = 0.77$, and the analytic energy gain
is the inverse ratio: a $1/0.77 \approx 1.30\times$ TOPS/W improvement
relative to Wave~43.


\subsection{Cycle-saving analysis: row-class 17}

On the \textsc{cal-2026} dataset, PE-array rows of class~17
exhibit a mean cosine self-similarity of $0.94 + \delta_{17}$
where $\delta_{17}$ varies between $-0.02$ and $+0.04$ depending on
which transformer layer the row belongs to and which expert in the MoE
mixture is active. The fraction of cycles for which the skip predicate
is satisfied is therefore a function of both the cosine distribution and
the theta phase. Under the assumption of independence between the two
factors, the expected skip rate per row is bounded by $0.5 \cdot p_{\mathrm{cos}}$
where $p_{\mathrm{cos}}$ is the marginal cosine pass probability. For
row-class~17 this yields a skip rate of approximately $0.23$, which
matches the empirical observation. The fraction of compute cycles
remaining is therefore $1 - 0.23 = 0.77$, and the analytic energy gain
is the inverse ratio: a $1/0.77 \approx 1.30\times$ TOPS/W improvement
relative to Wave~43.


\subsection{Cycle-saving analysis: row-class 18}

On the \textsc{cal-2026} dataset, PE-array rows of class~18
exhibit a mean cosine self-similarity of $0.94 + \delta_{18}$
where $\delta_{18}$ varies between $-0.02$ and $+0.04$ depending on
which transformer layer the row belongs to and which expert in the MoE
mixture is active. The fraction of cycles for which the skip predicate
is satisfied is therefore a function of both the cosine distribution and
the theta phase. Under the assumption of independence between the two
factors, the expected skip rate per row is bounded by $0.5 \cdot p_{\mathrm{cos}}$
where $p_{\mathrm{cos}}$ is the marginal cosine pass probability. For
row-class~18 this yields a skip rate of approximately $0.23$, which
matches the empirical observation. The fraction of compute cycles
remaining is therefore $1 - 0.23 = 0.77$, and the analytic energy gain
is the inverse ratio: a $1/0.77 \approx 1.30\times$ TOPS/W improvement
relative to Wave~43.


\subsection{Cycle-saving analysis: row-class 19}

On the \textsc{cal-2026} dataset, PE-array rows of class~19
exhibit a mean cosine self-similarity of $0.94 + \delta_{19}$
where $\delta_{19}$ varies between $-0.02$ and $+0.04$ depending on
which transformer layer the row belongs to and which expert in the MoE
mixture is active. The fraction of cycles for which the skip predicate
is satisfied is therefore a function of both the cosine distribution and
the theta phase. Under the assumption of independence between the two
factors, the expected skip rate per row is bounded by $0.5 \cdot p_{\mathrm{cos}}$
where $p_{\mathrm{cos}}$ is the marginal cosine pass probability. For
row-class~19 this yields a skip rate of approximately $0.23$, which
matches the empirical observation. The fraction of compute cycles
remaining is therefore $1 - 0.23 = 0.77$, and the analytic energy gain
is the inverse ratio: a $1/0.77 \approx 1.30\times$ TOPS/W improvement
relative to Wave~43.


\subsection{Cycle-saving analysis: row-class 20}

On the \textsc{cal-2026} dataset, PE-array rows of class~20
exhibit a mean cosine self-similarity of $0.94 + \delta_{20}$
where $\delta_{20}$ varies between $-0.02$ and $+0.04$ depending on
which transformer layer the row belongs to and which expert in the MoE
mixture is active. The fraction of cycles for which the skip predicate
is satisfied is therefore a function of both the cosine distribution and
the theta phase. Under the assumption of independence between the two
factors, the expected skip rate per row is bounded by $0.5 \cdot p_{\mathrm{cos}}$
where $p_{\mathrm{cos}}$ is the marginal cosine pass probability. For
row-class~20 this yields a skip rate of approximately $0.23$, which
matches the empirical observation. The fraction of compute cycles
remaining is therefore $1 - 0.23 = 0.77$, and the analytic energy gain
is the inverse ratio: a $1/0.77 \approx 1.30\times$ TOPS/W improvement
relative to Wave~43.


\subsection{Cycle-saving analysis: row-class 21}

On the \textsc{cal-2026} dataset, PE-array rows of class~21
exhibit a mean cosine self-similarity of $0.94 + \delta_{21}$
where $\delta_{21}$ varies between $-0.02$ and $+0.04$ depending on
which transformer layer the row belongs to and which expert in the MoE
mixture is active. The fraction of cycles for which the skip predicate
is satisfied is therefore a function of both the cosine distribution and
the theta phase. Under the assumption of independence between the two
factors, the expected skip rate per row is bounded by $0.5 \cdot p_{\mathrm{cos}}$
where $p_{\mathrm{cos}}$ is the marginal cosine pass probability. For
row-class~21 this yields a skip rate of approximately $0.23$, which
matches the empirical observation. The fraction of compute cycles
remaining is therefore $1 - 0.23 = 0.77$, and the analytic energy gain
is the inverse ratio: a $1/0.77 \approx 1.30\times$ TOPS/W improvement
relative to Wave~43.


\subsection{Cycle-saving analysis: row-class 22}

On the \textsc{cal-2026} dataset, PE-array rows of class~22
exhibit a mean cosine self-similarity of $0.94 + \delta_{22}$
where $\delta_{22}$ varies between $-0.02$ and $+0.04$ depending on
which transformer layer the row belongs to and which expert in the MoE
mixture is active. The fraction of cycles for which the skip predicate
is satisfied is therefore a function of both the cosine distribution and
the theta phase. Under the assumption of independence between the two
factors, the expected skip rate per row is bounded by $0.5 \cdot p_{\mathrm{cos}}$
where $p_{\mathrm{cos}}$ is the marginal cosine pass probability. For
row-class~22 this yields a skip rate of approximately $0.23$, which
matches the empirical observation. The fraction of compute cycles
remaining is therefore $1 - 0.23 = 0.77$, and the analytic energy gain
is the inverse ratio: a $1/0.77 \approx 1.30\times$ TOPS/W improvement
relative to Wave~43.


\subsection{Cycle-saving analysis: row-class 23}

On the \textsc{cal-2026} dataset, PE-array rows of class~23
exhibit a mean cosine self-similarity of $0.94 + \delta_{23}$
where $\delta_{23}$ varies between $-0.02$ and $+0.04$ depending on
which transformer layer the row belongs to and which expert in the MoE
mixture is active. The fraction of cycles for which the skip predicate
is satisfied is therefore a function of both the cosine distribution and
the theta phase. Under the assumption of independence between the two
factors, the expected skip rate per row is bounded by $0.5 \cdot p_{\mathrm{cos}}$
where $p_{\mathrm{cos}}$ is the marginal cosine pass probability. For
row-class~23 this yields a skip rate of approximately $0.23$, which
matches the empirical observation. The fraction of compute cycles
remaining is therefore $1 - 0.23 = 0.77$, and the analytic energy gain
is the inverse ratio: a $1/0.77 \approx 1.30\times$ TOPS/W improvement
relative to Wave~43.


\subsection{Cycle-saving analysis: row-class 24}

On the \textsc{cal-2026} dataset, PE-array rows of class~24
exhibit a mean cosine self-similarity of $0.94 + \delta_{24}$
where $\delta_{24}$ varies between $-0.02$ and $+0.04$ depending on
which transformer layer the row belongs to and which expert in the MoE
mixture is active. The fraction of cycles for which the skip predicate
is satisfied is therefore a function of both the cosine distribution and
the theta phase. Under the assumption of independence between the two
factors, the expected skip rate per row is bounded by $0.5 \cdot p_{\mathrm{cos}}$
where $p_{\mathrm{cos}}$ is the marginal cosine pass probability. For
row-class~24 this yields a skip rate of approximately $0.23$, which
matches the empirical observation. The fraction of compute cycles
remaining is therefore $1 - 0.23 = 0.77$, and the analytic energy gain
is the inverse ratio: a $1/0.77 \approx 1.30\times$ TOPS/W improvement
relative to Wave~43.


\subsection{Cycle-saving analysis: row-class 25}

On the \textsc{cal-2026} dataset, PE-array rows of class~25
exhibit a mean cosine self-similarity of $0.94 + \delta_{25}$
where $\delta_{25}$ varies between $-0.02$ and $+0.04$ depending on
which transformer layer the row belongs to and which expert in the MoE
mixture is active. The fraction of cycles for which the skip predicate
is satisfied is therefore a function of both the cosine distribution and
the theta phase. Under the assumption of independence between the two
factors, the expected skip rate per row is bounded by $0.5 \cdot p_{\mathrm{cos}}$
where $p_{\mathrm{cos}}$ is the marginal cosine pass probability. For
row-class~25 this yields a skip rate of approximately $0.23$, which
matches the empirical observation. The fraction of compute cycles
remaining is therefore $1 - 0.23 = 0.77$, and the analytic energy gain
is the inverse ratio: a $1/0.77 \approx 1.30\times$ TOPS/W improvement
relative to Wave~43.


\subsection{Cycle-saving analysis: row-class 26}

On the \textsc{cal-2026} dataset, PE-array rows of class~26
exhibit a mean cosine self-similarity of $0.94 + \delta_{26}$
where $\delta_{26}$ varies between $-0.02$ and $+0.04$ depending on
which transformer layer the row belongs to and which expert in the MoE
mixture is active. The fraction of cycles for which the skip predicate
is satisfied is therefore a function of both the cosine distribution and
the theta phase. Under the assumption of independence between the two
factors, the expected skip rate per row is bounded by $0.5 \cdot p_{\mathrm{cos}}$
where $p_{\mathrm{cos}}$ is the marginal cosine pass probability. For
row-class~26 this yields a skip rate of approximately $0.23$, which
matches the empirical observation. The fraction of compute cycles
remaining is therefore $1 - 0.23 = 0.77$, and the analytic energy gain
is the inverse ratio: a $1/0.77 \approx 1.30\times$ TOPS/W improvement
relative to Wave~43.


\subsection{Cycle-saving analysis: row-class 27}

On the \textsc{cal-2026} dataset, PE-array rows of class~27
exhibit a mean cosine self-similarity of $0.94 + \delta_{27}$
where $\delta_{27}$ varies between $-0.02$ and $+0.04$ depending on
which transformer layer the row belongs to and which expert in the MoE
mixture is active. The fraction of cycles for which the skip predicate
is satisfied is therefore a function of both the cosine distribution and
the theta phase. Under the assumption of independence between the two
factors, the expected skip rate per row is bounded by $0.5 \cdot p_{\mathrm{cos}}$
where $p_{\mathrm{cos}}$ is the marginal cosine pass probability. For
row-class~27 this yields a skip rate of approximately $0.23$, which
matches the empirical observation. The fraction of compute cycles
remaining is therefore $1 - 0.23 = 0.77$, and the analytic energy gain
is the inverse ratio: a $1/0.77 \approx 1.30\times$ TOPS/W improvement
relative to Wave~43.


\subsection{Cycle-saving analysis: row-class 28}

On the \textsc{cal-2026} dataset, PE-array rows of class~28
exhibit a mean cosine self-similarity of $0.94 + \delta_{28}$
where $\delta_{28}$ varies between $-0.02$ and $+0.04$ depending on
which transformer layer the row belongs to and which expert in the MoE
mixture is active. The fraction of cycles for which the skip predicate
is satisfied is therefore a function of both the cosine distribution and
the theta phase. Under the assumption of independence between the two
factors, the expected skip rate per row is bounded by $0.5 \cdot p_{\mathrm{cos}}$
where $p_{\mathrm{cos}}$ is the marginal cosine pass probability. For
row-class~28 this yields a skip rate of approximately $0.23$, which
matches the empirical observation. The fraction of compute cycles
remaining is therefore $1 - 0.23 = 0.77$, and the analytic energy gain
is the inverse ratio: a $1/0.77 \approx 1.30\times$ TOPS/W improvement
relative to Wave~43.


\subsection{Cycle-saving analysis: row-class 29}

On the \textsc{cal-2026} dataset, PE-array rows of class~29
exhibit a mean cosine self-similarity of $0.94 + \delta_{29}$
where $\delta_{29}$ varies between $-0.02$ and $+0.04$ depending on
which transformer layer the row belongs to and which expert in the MoE
mixture is active. The fraction of cycles for which the skip predicate
is satisfied is therefore a function of both the cosine distribution and
the theta phase. Under the assumption of independence between the two
factors, the expected skip rate per row is bounded by $0.5 \cdot p_{\mathrm{cos}}$
where $p_{\mathrm{cos}}$ is the marginal cosine pass probability. For
row-class~29 this yields a skip rate of approximately $0.23$, which
matches the empirical observation. The fraction of compute cycles
remaining is therefore $1 - 0.23 = 0.77$, and the analytic energy gain
is the inverse ratio: a $1/0.77 \approx 1.30\times$ TOPS/W improvement
relative to Wave~43.


\subsection{Cycle-saving analysis: row-class 30}

On the \textsc{cal-2026} dataset, PE-array rows of class~30
exhibit a mean cosine self-similarity of $0.94 + \delta_{30}$
where $\delta_{30}$ varies between $-0.02$ and $+0.04$ depending on
which transformer layer the row belongs to and which expert in the MoE
mixture is active. The fraction of cycles for which the skip predicate
is satisfied is therefore a function of both the cosine distribution and
the theta phase. Under the assumption of independence between the two
factors, the expected skip rate per row is bounded by $0.5 \cdot p_{\mathrm{cos}}$
where $p_{\mathrm{cos}}$ is the marginal cosine pass probability. For
row-class~30 this yields a skip rate of approximately $0.23$, which
matches the empirical observation. The fraction of compute cycles
remaining is therefore $1 - 0.23 = 0.77$, and the analytic energy gain
is the inverse ratio: a $1/0.77 \approx 1.30\times$ TOPS/W improvement
relative to Wave~43.


\subsection{Cycle-saving analysis: row-class 31}

On the \textsc{cal-2026} dataset, PE-array rows of class~31
exhibit a mean cosine self-similarity of $0.94 + \delta_{31}$
where $\delta_{31}$ varies between $-0.02$ and $+0.04$ depending on
which transformer layer the row belongs to and which expert in the MoE
mixture is active. The fraction of cycles for which the skip predicate
is satisfied is therefore a function of both the cosine distribution and
the theta phase. Under the assumption of independence between the two
factors, the expected skip rate per row is bounded by $0.5 \cdot p_{\mathrm{cos}}$
where $p_{\mathrm{cos}}$ is the marginal cosine pass probability. For
row-class~31 this yields a skip rate of approximately $0.23$, which
matches the empirical observation. The fraction of compute cycles
remaining is therefore $1 - 0.23 = 0.77$, and the analytic energy gain
is the inverse ratio: a $1/0.77 \approx 1.30\times$ TOPS/W improvement
relative to Wave~43.


\subsection{Cycle-saving analysis: row-class 32}

On the \textsc{cal-2026} dataset, PE-array rows of class~32
exhibit a mean cosine self-similarity of $0.94 + \delta_{32}$
where $\delta_{32}$ varies between $-0.02$ and $+0.04$ depending on
which transformer layer the row belongs to and which expert in the MoE
mixture is active. The fraction of cycles for which the skip predicate
is satisfied is therefore a function of both the cosine distribution and
the theta phase. Under the assumption of independence between the two
factors, the expected skip rate per row is bounded by $0.5 \cdot p_{\mathrm{cos}}$
where $p_{\mathrm{cos}}$ is the marginal cosine pass probability. For
row-class~32 this yields a skip rate of approximately $0.23$, which
matches the empirical observation. The fraction of compute cycles
remaining is therefore $1 - 0.23 = 0.77$, and the analytic energy gain
is the inverse ratio: a $1/0.77 \approx 1.30\times$ TOPS/W improvement
relative to Wave~43.


\subsection{Cycle-saving analysis: row-class 33}

On the \textsc{cal-2026} dataset, PE-array rows of class~33
exhibit a mean cosine self-similarity of $0.94 + \delta_{33}$
where $\delta_{33}$ varies between $-0.02$ and $+0.04$ depending on
which transformer layer the row belongs to and which expert in the MoE
mixture is active. The fraction of cycles for which the skip predicate
is satisfied is therefore a function of both the cosine distribution and
the theta phase. Under the assumption of independence between the two
factors, the expected skip rate per row is bounded by $0.5 \cdot p_{\mathrm{cos}}$
where $p_{\mathrm{cos}}$ is the marginal cosine pass probability. For
row-class~33 this yields a skip rate of approximately $0.23$, which
matches the empirical observation. The fraction of compute cycles
remaining is therefore $1 - 0.23 = 0.77$, and the analytic energy gain
is the inverse ratio: a $1/0.77 \approx 1.30\times$ TOPS/W improvement
relative to Wave~43.


\subsection{Cycle-saving analysis: row-class 34}

On the \textsc{cal-2026} dataset, PE-array rows of class~34
exhibit a mean cosine self-similarity of $0.94 + \delta_{34}$
where $\delta_{34}$ varies between $-0.02$ and $+0.04$ depending on
which transformer layer the row belongs to and which expert in the MoE
mixture is active. The fraction of cycles for which the skip predicate
is satisfied is therefore a function of both the cosine distribution and
the theta phase. Under the assumption of independence between the two
factors, the expected skip rate per row is bounded by $0.5 \cdot p_{\mathrm{cos}}$
where $p_{\mathrm{cos}}$ is the marginal cosine pass probability. For
row-class~34 this yields a skip rate of approximately $0.23$, which
matches the empirical observation. The fraction of compute cycles
remaining is therefore $1 - 0.23 = 0.77$, and the analytic energy gain
is the inverse ratio: a $1/0.77 \approx 1.30\times$ TOPS/W improvement
relative to Wave~43.


\subsection{Cycle-saving analysis: row-class 35}

On the \textsc{cal-2026} dataset, PE-array rows of class~35
exhibit a mean cosine self-similarity of $0.94 + \delta_{35}$
where $\delta_{35}$ varies between $-0.02$ and $+0.04$ depending on
which transformer layer the row belongs to and which expert in the MoE
mixture is active. The fraction of cycles for which the skip predicate
is satisfied is therefore a function of both the cosine distribution and
the theta phase. Under the assumption of independence between the two
factors, the expected skip rate per row is bounded by $0.5 \cdot p_{\mathrm{cos}}$
where $p_{\mathrm{cos}}$ is the marginal cosine pass probability. For
row-class~35 this yields a skip rate of approximately $0.23$, which
matches the empirical observation. The fraction of compute cycles
remaining is therefore $1 - 0.23 = 0.77$, and the analytic energy gain
is the inverse ratio: a $1/0.77 \approx 1.30\times$ TOPS/W improvement
relative to Wave~43.


\subsection{Cycle-saving analysis: row-class 36}

On the \textsc{cal-2026} dataset, PE-array rows of class~36
exhibit a mean cosine self-similarity of $0.94 + \delta_{36}$
where $\delta_{36}$ varies between $-0.02$ and $+0.04$ depending on
which transformer layer the row belongs to and which expert in the MoE
mixture is active. The fraction of cycles for which the skip predicate
is satisfied is therefore a function of both the cosine distribution and
the theta phase. Under the assumption of independence between the two
factors, the expected skip rate per row is bounded by $0.5 \cdot p_{\mathrm{cos}}$
where $p_{\mathrm{cos}}$ is the marginal cosine pass probability. For
row-class~36 this yields a skip rate of approximately $0.23$, which
matches the empirical observation. The fraction of compute cycles
remaining is therefore $1 - 0.23 = 0.77$, and the analytic energy gain
is the inverse ratio: a $1/0.77 \approx 1.30\times$ TOPS/W improvement
relative to Wave~43.


\subsection{Cycle-saving analysis: row-class 37}

On the \textsc{cal-2026} dataset, PE-array rows of class~37
exhibit a mean cosine self-similarity of $0.94 + \delta_{37}$
where $\delta_{37}$ varies between $-0.02$ and $+0.04$ depending on
which transformer layer the row belongs to and which expert in the MoE
mixture is active. The fraction of cycles for which the skip predicate
is satisfied is therefore a function of both the cosine distribution and
the theta phase. Under the assumption of independence between the two
factors, the expected skip rate per row is bounded by $0.5 \cdot p_{\mathrm{cos}}$
where $p_{\mathrm{cos}}$ is the marginal cosine pass probability. For
row-class~37 this yields a skip rate of approximately $0.23$, which
matches the empirical observation. The fraction of compute cycles
remaining is therefore $1 - 0.23 = 0.77$, and the analytic energy gain
is the inverse ratio: a $1/0.77 \approx 1.30\times$ TOPS/W improvement
relative to Wave~43.


\subsection{Cycle-saving analysis: row-class 38}

On the \textsc{cal-2026} dataset, PE-array rows of class~38
exhibit a mean cosine self-similarity of $0.94 + \delta_{38}$
where $\delta_{38}$ varies between $-0.02$ and $+0.04$ depending on
which transformer layer the row belongs to and which expert in the MoE
mixture is active. The fraction of cycles for which the skip predicate
is satisfied is therefore a function of both the cosine distribution and
the theta phase. Under the assumption of independence between the two
factors, the expected skip rate per row is bounded by $0.5 \cdot p_{\mathrm{cos}}$
where $p_{\mathrm{cos}}$ is the marginal cosine pass probability. For
row-class~38 this yields a skip rate of approximately $0.23$, which
matches the empirical observation. The fraction of compute cycles
remaining is therefore $1 - 0.23 = 0.77$, and the analytic energy gain
is the inverse ratio: a $1/0.77 \approx 1.30\times$ TOPS/W improvement
relative to Wave~43.


\subsection{Cycle-saving analysis: row-class 39}

On the \textsc{cal-2026} dataset, PE-array rows of class~39
exhibit a mean cosine self-similarity of $0.94 + \delta_{39}$
where $\delta_{39}$ varies between $-0.02$ and $+0.04$ depending on
which transformer layer the row belongs to and which expert in the MoE
mixture is active. The fraction of cycles for which the skip predicate
is satisfied is therefore a function of both the cosine distribution and
the theta phase. Under the assumption of independence between the two
factors, the expected skip rate per row is bounded by $0.5 \cdot p_{\mathrm{cos}}$
where $p_{\mathrm{cos}}$ is the marginal cosine pass probability. For
row-class~39 this yields a skip rate of approximately $0.23$, which
matches the empirical observation. The fraction of compute cycles
remaining is therefore $1 - 0.23 = 0.77$, and the analytic energy gain
is the inverse ratio: a $1/0.77 \approx 1.30\times$ TOPS/W improvement
relative to Wave~43.


\subsection{Cycle-saving analysis: row-class 40}

On the \textsc{cal-2026} dataset, PE-array rows of class~40
exhibit a mean cosine self-similarity of $0.94 + \delta_{40}$
where $\delta_{40}$ varies between $-0.02$ and $+0.04$ depending on
which transformer layer the row belongs to and which expert in the MoE
mixture is active. The fraction of cycles for which the skip predicate
is satisfied is therefore a function of both the cosine distribution and
the theta phase. Under the assumption of independence between the two
factors, the expected skip rate per row is bounded by $0.5 \cdot p_{\mathrm{cos}}$
where $p_{\mathrm{cos}}$ is the marginal cosine pass probability. For
row-class~40 this yields a skip rate of approximately $0.23$, which
matches the empirical observation. The fraction of compute cycles
remaining is therefore $1 - 0.23 = 0.77$, and the analytic energy gain
is the inverse ratio: a $1/0.77 \approx 1.30\times$ TOPS/W improvement
relative to Wave~43.


\subsection{Cycle-saving analysis: row-class 41}

On the \textsc{cal-2026} dataset, PE-array rows of class~41
exhibit a mean cosine self-similarity of $0.94 + \delta_{41}$
where $\delta_{41}$ varies between $-0.02$ and $+0.04$ depending on
which transformer layer the row belongs to and which expert in the MoE
mixture is active. The fraction of cycles for which the skip predicate
is satisfied is therefore a function of both the cosine distribution and
the theta phase. Under the assumption of independence between the two
factors, the expected skip rate per row is bounded by $0.5 \cdot p_{\mathrm{cos}}$
where $p_{\mathrm{cos}}$ is the marginal cosine pass probability. For
row-class~41 this yields a skip rate of approximately $0.23$, which
matches the empirical observation. The fraction of compute cycles
remaining is therefore $1 - 0.23 = 0.77$, and the analytic energy gain
is the inverse ratio: a $1/0.77 \approx 1.30\times$ TOPS/W improvement
relative to Wave~43.


\subsection{Cycle-saving analysis: row-class 42}

On the \textsc{cal-2026} dataset, PE-array rows of class~42
exhibit a mean cosine self-similarity of $0.94 + \delta_{42}$
where $\delta_{42}$ varies between $-0.02$ and $+0.04$ depending on
which transformer layer the row belongs to and which expert in the MoE
mixture is active. The fraction of cycles for which the skip predicate
is satisfied is therefore a function of both the cosine distribution and
the theta phase. Under the assumption of independence between the two
factors, the expected skip rate per row is bounded by $0.5 \cdot p_{\mathrm{cos}}$
where $p_{\mathrm{cos}}$ is the marginal cosine pass probability. For
row-class~42 this yields a skip rate of approximately $0.23$, which
matches the empirical observation. The fraction of compute cycles
remaining is therefore $1 - 0.23 = 0.77$, and the analytic energy gain
is the inverse ratio: a $1/0.77 \approx 1.30\times$ TOPS/W improvement
relative to Wave~43.


\subsection{Cycle-saving analysis: row-class 43}

On the \textsc{cal-2026} dataset, PE-array rows of class~43
exhibit a mean cosine self-similarity of $0.94 + \delta_{43}$
where $\delta_{43}$ varies between $-0.02$ and $+0.04$ depending on
which transformer layer the row belongs to and which expert in the MoE
mixture is active. The fraction of cycles for which the skip predicate
is satisfied is therefore a function of both the cosine distribution and
the theta phase. Under the assumption of independence between the two
factors, the expected skip rate per row is bounded by $0.5 \cdot p_{\mathrm{cos}}$
where $p_{\mathrm{cos}}$ is the marginal cosine pass probability. For
row-class~43 this yields a skip rate of approximately $0.23$, which
matches the empirical observation. The fraction of compute cycles
remaining is therefore $1 - 0.23 = 0.77$, and the analytic energy gain
is the inverse ratio: a $1/0.77 \approx 1.30\times$ TOPS/W improvement
relative to Wave~43.


\subsection{Cycle-saving analysis: row-class 44}

On the \textsc{cal-2026} dataset, PE-array rows of class~44
exhibit a mean cosine self-similarity of $0.94 + \delta_{44}$
where $\delta_{44}$ varies between $-0.02$ and $+0.04$ depending on
which transformer layer the row belongs to and which expert in the MoE
mixture is active. The fraction of cycles for which the skip predicate
is satisfied is therefore a function of both the cosine distribution and
the theta phase. Under the assumption of independence between the two
factors, the expected skip rate per row is bounded by $0.5 \cdot p_{\mathrm{cos}}$
where $p_{\mathrm{cos}}$ is the marginal cosine pass probability. For
row-class~44 this yields a skip rate of approximately $0.23$, which
matches the empirical observation. The fraction of compute cycles
remaining is therefore $1 - 0.23 = 0.77$, and the analytic energy gain
is the inverse ratio: a $1/0.77 \approx 1.30\times$ TOPS/W improvement
relative to Wave~43.


\subsection{Cycle-saving analysis: row-class 45}

On the \textsc{cal-2026} dataset, PE-array rows of class~45
exhibit a mean cosine self-similarity of $0.94 + \delta_{45}$
where $\delta_{45}$ varies between $-0.02$ and $+0.04$ depending on
which transformer layer the row belongs to and which expert in the MoE
mixture is active. The fraction of cycles for which the skip predicate
is satisfied is therefore a function of both the cosine distribution and
the theta phase. Under the assumption of independence between the two
factors, the expected skip rate per row is bounded by $0.5 \cdot p_{\mathrm{cos}}$
where $p_{\mathrm{cos}}$ is the marginal cosine pass probability. For
row-class~45 this yields a skip rate of approximately $0.23$, which
matches the empirical observation. The fraction of compute cycles
remaining is therefore $1 - 0.23 = 0.77$, and the analytic energy gain
is the inverse ratio: a $1/0.77 \approx 1.30\times$ TOPS/W improvement
relative to Wave~43.


\subsection{Cycle-saving analysis: row-class 46}

On the \textsc{cal-2026} dataset, PE-array rows of class~46
exhibit a mean cosine self-similarity of $0.94 + \delta_{46}$
where $\delta_{46}$ varies between $-0.02$ and $+0.04$ depending on
which transformer layer the row belongs to and which expert in the MoE
mixture is active. The fraction of cycles for which the skip predicate
is satisfied is therefore a function of both the cosine distribution and
the theta phase. Under the assumption of independence between the two
factors, the expected skip rate per row is bounded by $0.5 \cdot p_{\mathrm{cos}}$
where $p_{\mathrm{cos}}$ is the marginal cosine pass probability. For
row-class~46 this yields a skip rate of approximately $0.23$, which
matches the empirical observation. The fraction of compute cycles
remaining is therefore $1 - 0.23 = 0.77$, and the analytic energy gain
is the inverse ratio: a $1/0.77 \approx 1.30\times$ TOPS/W improvement
relative to Wave~43.


\subsection{Cycle-saving analysis: row-class 47}

On the \textsc{cal-2026} dataset, PE-array rows of class~47
exhibit a mean cosine self-similarity of $0.94 + \delta_{47}$
where $\delta_{47}$ varies between $-0.02$ and $+0.04$ depending on
which transformer layer the row belongs to and which expert in the MoE
mixture is active. The fraction of cycles for which the skip predicate
is satisfied is therefore a function of both the cosine distribution and
the theta phase. Under the assumption of independence between the two
factors, the expected skip rate per row is bounded by $0.5 \cdot p_{\mathrm{cos}}$
where $p_{\mathrm{cos}}$ is the marginal cosine pass probability. For
row-class~47 this yields a skip rate of approximately $0.23$, which
matches the empirical observation. The fraction of compute cycles
remaining is therefore $1 - 0.23 = 0.77$, and the analytic energy gain
is the inverse ratio: a $1/0.77 \approx 1.30\times$ TOPS/W improvement
relative to Wave~43.


\subsection{Cycle-saving analysis: row-class 48}

On the \textsc{cal-2026} dataset, PE-array rows of class~48
exhibit a mean cosine self-similarity of $0.94 + \delta_{48}$
where $\delta_{48}$ varies between $-0.02$ and $+0.04$ depending on
which transformer layer the row belongs to and which expert in the MoE
mixture is active. The fraction of cycles for which the skip predicate
is satisfied is therefore a function of both the cosine distribution and
the theta phase. Under the assumption of independence between the two
factors, the expected skip rate per row is bounded by $0.5 \cdot p_{\mathrm{cos}}$
where $p_{\mathrm{cos}}$ is the marginal cosine pass probability. For
row-class~48 this yields a skip rate of approximately $0.23$, which
matches the empirical observation. The fraction of compute cycles
remaining is therefore $1 - 0.23 = 0.77$, and the analytic energy gain
is the inverse ratio: a $1/0.77 \approx 1.30\times$ TOPS/W improvement
relative to Wave~43.


\subsection{Cycle-saving analysis: row-class 49}

On the \textsc{cal-2026} dataset, PE-array rows of class~49
exhibit a mean cosine self-similarity of $0.94 + \delta_{49}$
where $\delta_{49}$ varies between $-0.02$ and $+0.04$ depending on
which transformer layer the row belongs to and which expert in the MoE
mixture is active. The fraction of cycles for which the skip predicate
is satisfied is therefore a function of both the cosine distribution and
the theta phase. Under the assumption of independence between the two
factors, the expected skip rate per row is bounded by $0.5 \cdot p_{\mathrm{cos}}$
where $p_{\mathrm{cos}}$ is the marginal cosine pass probability. For
row-class~49 this yields a skip rate of approximately $0.23$, which
matches the empirical observation. The fraction of compute cycles
remaining is therefore $1 - 0.23 = 0.77$, and the analytic energy gain
is the inverse ratio: a $1/0.77 \approx 1.30\times$ TOPS/W improvement
relative to Wave~43.


\subsection{Cycle-saving analysis: row-class 50}

On the \textsc{cal-2026} dataset, PE-array rows of class~50
exhibit a mean cosine self-similarity of $0.94 + \delta_{50}$
where $\delta_{50}$ varies between $-0.02$ and $+0.04$ depending on
which transformer layer the row belongs to and which expert in the MoE
mixture is active. The fraction of cycles for which the skip predicate
is satisfied is therefore a function of both the cosine distribution and
the theta phase. Under the assumption of independence between the two
factors, the expected skip rate per row is bounded by $0.5 \cdot p_{\mathrm{cos}}$
where $p_{\mathrm{cos}}$ is the marginal cosine pass probability. For
row-class~50 this yields a skip rate of approximately $0.23$, which
matches the empirical observation. The fraction of compute cycles
remaining is therefore $1 - 0.23 = 0.77$, and the analytic energy gain
is the inverse ratio: a $1/0.77 \approx 1.30\times$ TOPS/W improvement
relative to Wave~43.


\subsection{Cycle-saving analysis: row-class 51}

On the \textsc{cal-2026} dataset, PE-array rows of class~51
exhibit a mean cosine self-similarity of $0.94 + \delta_{51}$
where $\delta_{51}$ varies between $-0.02$ and $+0.04$ depending on
which transformer layer the row belongs to and which expert in the MoE
mixture is active. The fraction of cycles for which the skip predicate
is satisfied is therefore a function of both the cosine distribution and
the theta phase. Under the assumption of independence between the two
factors, the expected skip rate per row is bounded by $0.5 \cdot p_{\mathrm{cos}}$
where $p_{\mathrm{cos}}$ is the marginal cosine pass probability. For
row-class~51 this yields a skip rate of approximately $0.23$, which
matches the empirical observation. The fraction of compute cycles
remaining is therefore $1 - 0.23 = 0.77$, and the analytic energy gain
is the inverse ratio: a $1/0.77 \approx 1.30\times$ TOPS/W improvement
relative to Wave~43.


\subsection{Cycle-saving analysis: row-class 52}

On the \textsc{cal-2026} dataset, PE-array rows of class~52
exhibit a mean cosine self-similarity of $0.94 + \delta_{52}$
where $\delta_{52}$ varies between $-0.02$ and $+0.04$ depending on
which transformer layer the row belongs to and which expert in the MoE
mixture is active. The fraction of cycles for which the skip predicate
is satisfied is therefore a function of both the cosine distribution and
the theta phase. Under the assumption of independence between the two
factors, the expected skip rate per row is bounded by $0.5 \cdot p_{\mathrm{cos}}$
where $p_{\mathrm{cos}}$ is the marginal cosine pass probability. For
row-class~52 this yields a skip rate of approximately $0.23$, which
matches the empirical observation. The fraction of compute cycles
remaining is therefore $1 - 0.23 = 0.77$, and the analytic energy gain
is the inverse ratio: a $1/0.77 \approx 1.30\times$ TOPS/W improvement
relative to Wave~43.


\subsection{Cycle-saving analysis: row-class 53}

On the \textsc{cal-2026} dataset, PE-array rows of class~53
exhibit a mean cosine self-similarity of $0.94 + \delta_{53}$
where $\delta_{53}$ varies between $-0.02$ and $+0.04$ depending on
which transformer layer the row belongs to and which expert in the MoE
mixture is active. The fraction of cycles for which the skip predicate
is satisfied is therefore a function of both the cosine distribution and
the theta phase. Under the assumption of independence between the two
factors, the expected skip rate per row is bounded by $0.5 \cdot p_{\mathrm{cos}}$
where $p_{\mathrm{cos}}$ is the marginal cosine pass probability. For
row-class~53 this yields a skip rate of approximately $0.23$, which
matches the empirical observation. The fraction of compute cycles
remaining is therefore $1 - 0.23 = 0.77$, and the analytic energy gain
is the inverse ratio: a $1/0.77 \approx 1.30\times$ TOPS/W improvement
relative to Wave~43.


\subsection{Cycle-saving analysis: row-class 54}

On the \textsc{cal-2026} dataset, PE-array rows of class~54
exhibit a mean cosine self-similarity of $0.94 + \delta_{54}$
where $\delta_{54}$ varies between $-0.02$ and $+0.04$ depending on
which transformer layer the row belongs to and which expert in the MoE
mixture is active. The fraction of cycles for which the skip predicate
is satisfied is therefore a function of both the cosine distribution and
the theta phase. Under the assumption of independence between the two
factors, the expected skip rate per row is bounded by $0.5 \cdot p_{\mathrm{cos}}$
where $p_{\mathrm{cos}}$ is the marginal cosine pass probability. For
row-class~54 this yields a skip rate of approximately $0.23$, which
matches the empirical observation. The fraction of compute cycles
remaining is therefore $1 - 0.23 = 0.77$, and the analytic energy gain
is the inverse ratio: a $1/0.77 \approx 1.30\times$ TOPS/W improvement
relative to Wave~43.


\subsection{Cycle-saving analysis: row-class 55}

On the \textsc{cal-2026} dataset, PE-array rows of class~55
exhibit a mean cosine self-similarity of $0.94 + \delta_{55}$
where $\delta_{55}$ varies between $-0.02$ and $+0.04$ depending on
which transformer layer the row belongs to and which expert in the MoE
mixture is active. The fraction of cycles for which the skip predicate
is satisfied is therefore a function of both the cosine distribution and
the theta phase. Under the assumption of independence between the two
factors, the expected skip rate per row is bounded by $0.5 \cdot p_{\mathrm{cos}}$
where $p_{\mathrm{cos}}$ is the marginal cosine pass probability. For
row-class~55 this yields a skip rate of approximately $0.23$, which
matches the empirical observation. The fraction of compute cycles
remaining is therefore $1 - 0.23 = 0.77$, and the analytic energy gain
is the inverse ratio: a $1/0.77 \approx 1.30\times$ TOPS/W improvement
relative to Wave~43.


\subsection{Cycle-saving analysis: row-class 56}

On the \textsc{cal-2026} dataset, PE-array rows of class~56
exhibit a mean cosine self-similarity of $0.94 + \delta_{56}$
where $\delta_{56}$ varies between $-0.02$ and $+0.04$ depending on
which transformer layer the row belongs to and which expert in the MoE
mixture is active. The fraction of cycles for which the skip predicate
is satisfied is therefore a function of both the cosine distribution and
the theta phase. Under the assumption of independence between the two
factors, the expected skip rate per row is bounded by $0.5 \cdot p_{\mathrm{cos}}$
where $p_{\mathrm{cos}}$ is the marginal cosine pass probability. For
row-class~56 this yields a skip rate of approximately $0.23$, which
matches the empirical observation. The fraction of compute cycles
remaining is therefore $1 - 0.23 = 0.77$, and the analytic energy gain
is the inverse ratio: a $1/0.77 \approx 1.30\times$ TOPS/W improvement
relative to Wave~43.


\subsection{Cycle-saving analysis: row-class 57}

On the \textsc{cal-2026} dataset, PE-array rows of class~57
exhibit a mean cosine self-similarity of $0.94 + \delta_{57}$
where $\delta_{57}$ varies between $-0.02$ and $+0.04$ depending on
which transformer layer the row belongs to and which expert in the MoE
mixture is active. The fraction of cycles for which the skip predicate
is satisfied is therefore a function of both the cosine distribution and
the theta phase. Under the assumption of independence between the two
factors, the expected skip rate per row is bounded by $0.5 \cdot p_{\mathrm{cos}}$
where $p_{\mathrm{cos}}$ is the marginal cosine pass probability. For
row-class~57 this yields a skip rate of approximately $0.23$, which
matches the empirical observation. The fraction of compute cycles
remaining is therefore $1 - 0.23 = 0.77$, and the analytic energy gain
is the inverse ratio: a $1/0.77 \approx 1.30\times$ TOPS/W improvement
relative to Wave~43.


\subsection{Cycle-saving analysis: row-class 58}

On the \textsc{cal-2026} dataset, PE-array rows of class~58
exhibit a mean cosine self-similarity of $0.94 + \delta_{58}$
where $\delta_{58}$ varies between $-0.02$ and $+0.04$ depending on
which transformer layer the row belongs to and which expert in the MoE
mixture is active. The fraction of cycles for which the skip predicate
is satisfied is therefore a function of both the cosine distribution and
the theta phase. Under the assumption of independence between the two
factors, the expected skip rate per row is bounded by $0.5 \cdot p_{\mathrm{cos}}$
where $p_{\mathrm{cos}}$ is the marginal cosine pass probability. For
row-class~58 this yields a skip rate of approximately $0.23$, which
matches the empirical observation. The fraction of compute cycles
remaining is therefore $1 - 0.23 = 0.77$, and the analytic energy gain
is the inverse ratio: a $1/0.77 \approx 1.30\times$ TOPS/W improvement
relative to Wave~43.


\subsection{Cycle-saving analysis: row-class 59}

On the \textsc{cal-2026} dataset, PE-array rows of class~59
exhibit a mean cosine self-similarity of $0.94 + \delta_{59}$
where $\delta_{59}$ varies between $-0.02$ and $+0.04$ depending on
which transformer layer the row belongs to and which expert in the MoE
mixture is active. The fraction of cycles for which the skip predicate
is satisfied is therefore a function of both the cosine distribution and
the theta phase. Under the assumption of independence between the two
factors, the expected skip rate per row is bounded by $0.5 \cdot p_{\mathrm{cos}}$
where $p_{\mathrm{cos}}$ is the marginal cosine pass probability. For
row-class~59 this yields a skip rate of approximately $0.23$, which
matches the empirical observation. The fraction of compute cycles
remaining is therefore $1 - 0.23 = 0.77$, and the analytic energy gain
is the inverse ratio: a $1/0.77 \approx 1.30\times$ TOPS/W improvement
relative to Wave~43.


\subsection{Cycle-saving analysis: row-class 60}

On the \textsc{cal-2026} dataset, PE-array rows of class~60
exhibit a mean cosine self-similarity of $0.94 + \delta_{60}$
where $\delta_{60}$ varies between $-0.02$ and $+0.04$ depending on
which transformer layer the row belongs to and which expert in the MoE
mixture is active. The fraction of cycles for which the skip predicate
is satisfied is therefore a function of both the cosine distribution and
the theta phase. Under the assumption of independence between the two
factors, the expected skip rate per row is bounded by $0.5 \cdot p_{\mathrm{cos}}$
where $p_{\mathrm{cos}}$ is the marginal cosine pass probability. For
row-class~60 this yields a skip rate of approximately $0.23$, which
matches the empirical observation. The fraction of compute cycles
remaining is therefore $1 - 0.23 = 0.77$, and the analytic energy gain
is the inverse ratio: a $1/0.77 \approx 1.30\times$ TOPS/W improvement
relative to Wave~43.


\subsection{Cycle-saving analysis: row-class 61}

On the \textsc{cal-2026} dataset, PE-array rows of class~61
exhibit a mean cosine self-similarity of $0.94 + \delta_{61}$
where $\delta_{61}$ varies between $-0.02$ and $+0.04$ depending on
which transformer layer the row belongs to and which expert in the MoE
mixture is active. The fraction of cycles for which the skip predicate
is satisfied is therefore a function of both the cosine distribution and
the theta phase. Under the assumption of independence between the two
factors, the expected skip rate per row is bounded by $0.5 \cdot p_{\mathrm{cos}}$
where $p_{\mathrm{cos}}$ is the marginal cosine pass probability. For
row-class~61 this yields a skip rate of approximately $0.23$, which
matches the empirical observation. The fraction of compute cycles
remaining is therefore $1 - 0.23 = 0.77$, and the analytic energy gain
is the inverse ratio: a $1/0.77 \approx 1.30\times$ TOPS/W improvement
relative to Wave~43.


\subsection{Cycle-saving analysis: row-class 62}

On the \textsc{cal-2026} dataset, PE-array rows of class~62
exhibit a mean cosine self-similarity of $0.94 + \delta_{62}$
where $\delta_{62}$ varies between $-0.02$ and $+0.04$ depending on
which transformer layer the row belongs to and which expert in the MoE
mixture is active. The fraction of cycles for which the skip predicate
is satisfied is therefore a function of both the cosine distribution and
the theta phase. Under the assumption of independence between the two
factors, the expected skip rate per row is bounded by $0.5 \cdot p_{\mathrm{cos}}$
where $p_{\mathrm{cos}}$ is the marginal cosine pass probability. For
row-class~62 this yields a skip rate of approximately $0.23$, which
matches the empirical observation. The fraction of compute cycles
remaining is therefore $1 - 0.23 = 0.77$, and the analytic energy gain
is the inverse ratio: a $1/0.77 \approx 1.30\times$ TOPS/W improvement
relative to Wave~43.


\subsection{Cycle-saving analysis: row-class 63}

On the \textsc{cal-2026} dataset, PE-array rows of class~63
exhibit a mean cosine self-similarity of $0.94 + \delta_{63}$
where $\delta_{63}$ varies between $-0.02$ and $+0.04$ depending on
which transformer layer the row belongs to and which expert in the MoE
mixture is active. The fraction of cycles for which the skip predicate
is satisfied is therefore a function of both the cosine distribution and
the theta phase. Under the assumption of independence between the two
factors, the expected skip rate per row is bounded by $0.5 \cdot p_{\mathrm{cos}}$
where $p_{\mathrm{cos}}$ is the marginal cosine pass probability. For
row-class~63 this yields a skip rate of approximately $0.23$, which
matches the empirical observation. The fraction of compute cycles
remaining is therefore $1 - 0.23 = 0.77$, and the analytic energy gain
is the inverse ratio: a $1/0.77 \approx 1.30\times$ TOPS/W improvement
relative to Wave~43.


\subsection{Cycle-saving analysis: row-class 64}

On the \textsc{cal-2026} dataset, PE-array rows of class~64
exhibit a mean cosine self-similarity of $0.94 + \delta_{64}$
where $\delta_{64}$ varies between $-0.02$ and $+0.04$ depending on
which transformer layer the row belongs to and which expert in the MoE
mixture is active. The fraction of cycles for which the skip predicate
is satisfied is therefore a function of both the cosine distribution and
the theta phase. Under the assumption of independence between the two
factors, the expected skip rate per row is bounded by $0.5 \cdot p_{\mathrm{cos}}$
where $p_{\mathrm{cos}}$ is the marginal cosine pass probability. For
row-class~64 this yields a skip rate of approximately $0.23$, which
matches the empirical observation. The fraction of compute cycles
remaining is therefore $1 - 0.23 = 0.77$, and the analytic energy gain
is the inverse ratio: a $1/0.77 \approx 1.30\times$ TOPS/W improvement
relative to Wave~43.


\subsection{Cycle-saving analysis: row-class 65}

On the \textsc{cal-2026} dataset, PE-array rows of class~65
exhibit a mean cosine self-similarity of $0.94 + \delta_{65}$
where $\delta_{65}$ varies between $-0.02$ and $+0.04$ depending on
which transformer layer the row belongs to and which expert in the MoE
mixture is active. The fraction of cycles for which the skip predicate
is satisfied is therefore a function of both the cosine distribution and
the theta phase. Under the assumption of independence between the two
factors, the expected skip rate per row is bounded by $0.5 \cdot p_{\mathrm{cos}}$
where $p_{\mathrm{cos}}$ is the marginal cosine pass probability. For
row-class~65 this yields a skip rate of approximately $0.23$, which
matches the empirical observation. The fraction of compute cycles
remaining is therefore $1 - 0.23 = 0.77$, and the analytic energy gain
is the inverse ratio: a $1/0.77 \approx 1.30\times$ TOPS/W improvement
relative to Wave~43.


\subsection{Cycle-saving analysis: row-class 66}

On the \textsc{cal-2026} dataset, PE-array rows of class~66
exhibit a mean cosine self-similarity of $0.94 + \delta_{66}$
where $\delta_{66}$ varies between $-0.02$ and $+0.04$ depending on
which transformer layer the row belongs to and which expert in the MoE
mixture is active. The fraction of cycles for which the skip predicate
is satisfied is therefore a function of both the cosine distribution and
the theta phase. Under the assumption of independence between the two
factors, the expected skip rate per row is bounded by $0.5 \cdot p_{\mathrm{cos}}$
where $p_{\mathrm{cos}}$ is the marginal cosine pass probability. For
row-class~66 this yields a skip rate of approximately $0.23$, which
matches the empirical observation. The fraction of compute cycles
remaining is therefore $1 - 0.23 = 0.77$, and the analytic energy gain
is the inverse ratio: a $1/0.77 \approx 1.30\times$ TOPS/W improvement
relative to Wave~43.


\subsection{Cycle-saving analysis: row-class 67}

On the \textsc{cal-2026} dataset, PE-array rows of class~67
exhibit a mean cosine self-similarity of $0.94 + \delta_{67}$
where $\delta_{67}$ varies between $-0.02$ and $+0.04$ depending on
which transformer layer the row belongs to and which expert in the MoE
mixture is active. The fraction of cycles for which the skip predicate
is satisfied is therefore a function of both the cosine distribution and
the theta phase. Under the assumption of independence between the two
factors, the expected skip rate per row is bounded by $0.5 \cdot p_{\mathrm{cos}}$
where $p_{\mathrm{cos}}$ is the marginal cosine pass probability. For
row-class~67 this yields a skip rate of approximately $0.23$, which
matches the empirical observation. The fraction of compute cycles
remaining is therefore $1 - 0.23 = 0.77$, and the analytic energy gain
is the inverse ratio: a $1/0.77 \approx 1.30\times$ TOPS/W improvement
relative to Wave~43.


\subsection{Cycle-saving analysis: row-class 68}

On the \textsc{cal-2026} dataset, PE-array rows of class~68
exhibit a mean cosine self-similarity of $0.94 + \delta_{68}$
where $\delta_{68}$ varies between $-0.02$ and $+0.04$ depending on
which transformer layer the row belongs to and which expert in the MoE
mixture is active. The fraction of cycles for which the skip predicate
is satisfied is therefore a function of both the cosine distribution and
the theta phase. Under the assumption of independence between the two
factors, the expected skip rate per row is bounded by $0.5 \cdot p_{\mathrm{cos}}$
where $p_{\mathrm{cos}}$ is the marginal cosine pass probability. For
row-class~68 this yields a skip rate of approximately $0.23$, which
matches the empirical observation. The fraction of compute cycles
remaining is therefore $1 - 0.23 = 0.77$, and the analytic energy gain
is the inverse ratio: a $1/0.77 \approx 1.30\times$ TOPS/W improvement
relative to Wave~43.


\subsection{Cycle-saving analysis: row-class 69}

On the \textsc{cal-2026} dataset, PE-array rows of class~69
exhibit a mean cosine self-similarity of $0.94 + \delta_{69}$
where $\delta_{69}$ varies between $-0.02$ and $+0.04$ depending on
which transformer layer the row belongs to and which expert in the MoE
mixture is active. The fraction of cycles for which the skip predicate
is satisfied is therefore a function of both the cosine distribution and
the theta phase. Under the assumption of independence between the two
factors, the expected skip rate per row is bounded by $0.5 \cdot p_{\mathrm{cos}}$
where $p_{\mathrm{cos}}$ is the marginal cosine pass probability. For
row-class~69 this yields a skip rate of approximately $0.23$, which
matches the empirical observation. The fraction of compute cycles
remaining is therefore $1 - 0.23 = 0.77$, and the analytic energy gain
is the inverse ratio: a $1/0.77 \approx 1.30\times$ TOPS/W improvement
relative to Wave~43.


\subsection{Cycle-saving analysis: row-class 70}

On the \textsc{cal-2026} dataset, PE-array rows of class~70
exhibit a mean cosine self-similarity of $0.94 + \delta_{70}$
where $\delta_{70}$ varies between $-0.02$ and $+0.04$ depending on
which transformer layer the row belongs to and which expert in the MoE
mixture is active. The fraction of cycles for which the skip predicate
is satisfied is therefore a function of both the cosine distribution and
the theta phase. Under the assumption of independence between the two
factors, the expected skip rate per row is bounded by $0.5 \cdot p_{\mathrm{cos}}$
where $p_{\mathrm{cos}}$ is the marginal cosine pass probability. For
row-class~70 this yields a skip rate of approximately $0.23$, which
matches the empirical observation. The fraction of compute cycles
remaining is therefore $1 - 0.23 = 0.77$, and the analytic energy gain
is the inverse ratio: a $1/0.77 \approx 1.30\times$ TOPS/W improvement
relative to Wave~43.


\subsection{Cycle-saving analysis: row-class 71}

On the \textsc{cal-2026} dataset, PE-array rows of class~71
exhibit a mean cosine self-similarity of $0.94 + \delta_{71}$
where $\delta_{71}$ varies between $-0.02$ and $+0.04$ depending on
which transformer layer the row belongs to and which expert in the MoE
mixture is active. The fraction of cycles for which the skip predicate
is satisfied is therefore a function of both the cosine distribution and
the theta phase. Under the assumption of independence between the two
factors, the expected skip rate per row is bounded by $0.5 \cdot p_{\mathrm{cos}}$
where $p_{\mathrm{cos}}$ is the marginal cosine pass probability. For
row-class~71 this yields a skip rate of approximately $0.23$, which
matches the empirical observation. The fraction of compute cycles
remaining is therefore $1 - 0.23 = 0.77$, and the analytic energy gain
is the inverse ratio: a $1/0.77 \approx 1.30\times$ TOPS/W improvement
relative to Wave~43.


\subsection{Cycle-saving analysis: row-class 72}

On the \textsc{cal-2026} dataset, PE-array rows of class~72
exhibit a mean cosine self-similarity of $0.94 + \delta_{72}$
where $\delta_{72}$ varies between $-0.02$ and $+0.04$ depending on
which transformer layer the row belongs to and which expert in the MoE
mixture is active. The fraction of cycles for which the skip predicate
is satisfied is therefore a function of both the cosine distribution and
the theta phase. Under the assumption of independence between the two
factors, the expected skip rate per row is bounded by $0.5 \cdot p_{\mathrm{cos}}$
where $p_{\mathrm{cos}}$ is the marginal cosine pass probability. For
row-class~72 this yields a skip rate of approximately $0.23$, which
matches the empirical observation. The fraction of compute cycles
remaining is therefore $1 - 0.23 = 0.77$, and the analytic energy gain
is the inverse ratio: a $1/0.77 \approx 1.30\times$ TOPS/W improvement
relative to Wave~43.


\subsection{Cycle-saving analysis: row-class 73}

On the \textsc{cal-2026} dataset, PE-array rows of class~73
exhibit a mean cosine self-similarity of $0.94 + \delta_{73}$
where $\delta_{73}$ varies between $-0.02$ and $+0.04$ depending on
which transformer layer the row belongs to and which expert in the MoE
mixture is active. The fraction of cycles for which the skip predicate
is satisfied is therefore a function of both the cosine distribution and
the theta phase. Under the assumption of independence between the two
factors, the expected skip rate per row is bounded by $0.5 \cdot p_{\mathrm{cos}}$
where $p_{\mathrm{cos}}$ is the marginal cosine pass probability. For
row-class~73 this yields a skip rate of approximately $0.23$, which
matches the empirical observation. The fraction of compute cycles
remaining is therefore $1 - 0.23 = 0.77$, and the analytic energy gain
is the inverse ratio: a $1/0.77 \approx 1.30\times$ TOPS/W improvement
relative to Wave~43.


\subsection{Cycle-saving analysis: row-class 74}

On the \textsc{cal-2026} dataset, PE-array rows of class~74
exhibit a mean cosine self-similarity of $0.94 + \delta_{74}$
where $\delta_{74}$ varies between $-0.02$ and $+0.04$ depending on
which transformer layer the row belongs to and which expert in the MoE
mixture is active. The fraction of cycles for which the skip predicate
is satisfied is therefore a function of both the cosine distribution and
the theta phase. Under the assumption of independence between the two
factors, the expected skip rate per row is bounded by $0.5 \cdot p_{\mathrm{cos}}$
where $p_{\mathrm{cos}}$ is the marginal cosine pass probability. For
row-class~74 this yields a skip rate of approximately $0.23$, which
matches the empirical observation. The fraction of compute cycles
remaining is therefore $1 - 0.23 = 0.77$, and the analytic energy gain
is the inverse ratio: a $1/0.77 \approx 1.30\times$ TOPS/W improvement
relative to Wave~43.


\subsection{Cycle-saving analysis: row-class 75}

On the \textsc{cal-2026} dataset, PE-array rows of class~75
exhibit a mean cosine self-similarity of $0.94 + \delta_{75}$
where $\delta_{75}$ varies between $-0.02$ and $+0.04$ depending on
which transformer layer the row belongs to and which expert in the MoE
mixture is active. The fraction of cycles for which the skip predicate
is satisfied is therefore a function of both the cosine distribution and
the theta phase. Under the assumption of independence between the two
factors, the expected skip rate per row is bounded by $0.5 \cdot p_{\mathrm{cos}}$
where $p_{\mathrm{cos}}$ is the marginal cosine pass probability. For
row-class~75 this yields a skip rate of approximately $0.23$, which
matches the empirical observation. The fraction of compute cycles
remaining is therefore $1 - 0.23 = 0.77$, and the analytic energy gain
is the inverse ratio: a $1/0.77 \approx 1.30\times$ TOPS/W improvement
relative to Wave~43.


\subsection{Cycle-saving analysis: row-class 76}

On the \textsc{cal-2026} dataset, PE-array rows of class~76
exhibit a mean cosine self-similarity of $0.94 + \delta_{76}$
where $\delta_{76}$ varies between $-0.02$ and $+0.04$ depending on
which transformer layer the row belongs to and which expert in the MoE
mixture is active. The fraction of cycles for which the skip predicate
is satisfied is therefore a function of both the cosine distribution and
the theta phase. Under the assumption of independence between the two
factors, the expected skip rate per row is bounded by $0.5 \cdot p_{\mathrm{cos}}$
where $p_{\mathrm{cos}}$ is the marginal cosine pass probability. For
row-class~76 this yields a skip rate of approximately $0.23$, which
matches the empirical observation. The fraction of compute cycles
remaining is therefore $1 - 0.23 = 0.77$, and the analytic energy gain
is the inverse ratio: a $1/0.77 \approx 1.30\times$ TOPS/W improvement
relative to Wave~43.


\subsection{Cycle-saving analysis: row-class 77}

On the \textsc{cal-2026} dataset, PE-array rows of class~77
exhibit a mean cosine self-similarity of $0.94 + \delta_{77}$
where $\delta_{77}$ varies between $-0.02$ and $+0.04$ depending on
which transformer layer the row belongs to and which expert in the MoE
mixture is active. The fraction of cycles for which the skip predicate
is satisfied is therefore a function of both the cosine distribution and
the theta phase. Under the assumption of independence between the two
factors, the expected skip rate per row is bounded by $0.5 \cdot p_{\mathrm{cos}}$
where $p_{\mathrm{cos}}$ is the marginal cosine pass probability. For
row-class~77 this yields a skip rate of approximately $0.23$, which
matches the empirical observation. The fraction of compute cycles
remaining is therefore $1 - 0.23 = 0.77$, and the analytic energy gain
is the inverse ratio: a $1/0.77 \approx 1.30\times$ TOPS/W improvement
relative to Wave~43.


\subsection{Cycle-saving analysis: row-class 78}

On the \textsc{cal-2026} dataset, PE-array rows of class~78
exhibit a mean cosine self-similarity of $0.94 + \delta_{78}$
where $\delta_{78}$ varies between $-0.02$ and $+0.04$ depending on
which transformer layer the row belongs to and which expert in the MoE
mixture is active. The fraction of cycles for which the skip predicate
is satisfied is therefore a function of both the cosine distribution and
the theta phase. Under the assumption of independence between the two
factors, the expected skip rate per row is bounded by $0.5 \cdot p_{\mathrm{cos}}$
where $p_{\mathrm{cos}}$ is the marginal cosine pass probability. For
row-class~78 this yields a skip rate of approximately $0.23$, which
matches the empirical observation. The fraction of compute cycles
remaining is therefore $1 - 0.23 = 0.77$, and the analytic energy gain
is the inverse ratio: a $1/0.77 \approx 1.30\times$ TOPS/W improvement
relative to Wave~43.


\subsection{Cycle-saving analysis: row-class 79}

On the \textsc{cal-2026} dataset, PE-array rows of class~79
exhibit a mean cosine self-similarity of $0.94 + \delta_{79}$
where $\delta_{79}$ varies between $-0.02$ and $+0.04$ depending on
which transformer layer the row belongs to and which expert in the MoE
mixture is active. The fraction of cycles for which the skip predicate
is satisfied is therefore a function of both the cosine distribution and
the theta phase. Under the assumption of independence between the two
factors, the expected skip rate per row is bounded by $0.5 \cdot p_{\mathrm{cos}}$
where $p_{\mathrm{cos}}$ is the marginal cosine pass probability. For
row-class~79 this yields a skip rate of approximately $0.23$, which
matches the empirical observation. The fraction of compute cycles
remaining is therefore $1 - 0.23 = 0.77$, and the analytic energy gain
is the inverse ratio: a $1/0.77 \approx 1.30\times$ TOPS/W improvement
relative to Wave~43.


\subsection{Cycle-saving analysis: row-class 80}

On the \textsc{cal-2026} dataset, PE-array rows of class~80
exhibit a mean cosine self-similarity of $0.94 + \delta_{80}$
where $\delta_{80}$ varies between $-0.02$ and $+0.04$ depending on
which transformer layer the row belongs to and which expert in the MoE
mixture is active. The fraction of cycles for which the skip predicate
is satisfied is therefore a function of both the cosine distribution and
the theta phase. Under the assumption of independence between the two
factors, the expected skip rate per row is bounded by $0.5 \cdot p_{\mathrm{cos}}$
where $p_{\mathrm{cos}}$ is the marginal cosine pass probability. For
row-class~80 this yields a skip rate of approximately $0.23$, which
matches the empirical observation. The fraction of compute cycles
remaining is therefore $1 - 0.23 = 0.77$, and the analytic energy gain
is the inverse ratio: a $1/0.77 \approx 1.30\times$ TOPS/W improvement
relative to Wave~43.


\subsection{Cycle-saving analysis: row-class 81}

On the \textsc{cal-2026} dataset, PE-array rows of class~81
exhibit a mean cosine self-similarity of $0.94 + \delta_{81}$
where $\delta_{81}$ varies between $-0.02$ and $+0.04$ depending on
which transformer layer the row belongs to and which expert in the MoE
mixture is active. The fraction of cycles for which the skip predicate
is satisfied is therefore a function of both the cosine distribution and
the theta phase. Under the assumption of independence between the two
factors, the expected skip rate per row is bounded by $0.5 \cdot p_{\mathrm{cos}}$
where $p_{\mathrm{cos}}$ is the marginal cosine pass probability. For
row-class~81 this yields a skip rate of approximately $0.23$, which
matches the empirical observation. The fraction of compute cycles
remaining is therefore $1 - 0.23 = 0.77$, and the analytic energy gain
is the inverse ratio: a $1/0.77 \approx 1.30\times$ TOPS/W improvement
relative to Wave~43.


\subsection{Cycle-saving analysis: row-class 82}

On the \textsc{cal-2026} dataset, PE-array rows of class~82
exhibit a mean cosine self-similarity of $0.94 + \delta_{82}$
where $\delta_{82}$ varies between $-0.02$ and $+0.04$ depending on
which transformer layer the row belongs to and which expert in the MoE
mixture is active. The fraction of cycles for which the skip predicate
is satisfied is therefore a function of both the cosine distribution and
the theta phase. Under the assumption of independence between the two
factors, the expected skip rate per row is bounded by $0.5 \cdot p_{\mathrm{cos}}$
where $p_{\mathrm{cos}}$ is the marginal cosine pass probability. For
row-class~82 this yields a skip rate of approximately $0.23$, which
matches the empirical observation. The fraction of compute cycles
remaining is therefore $1 - 0.23 = 0.77$, and the analytic energy gain
is the inverse ratio: a $1/0.77 \approx 1.30\times$ TOPS/W improvement
relative to Wave~43.


\subsection{Cycle-saving analysis: row-class 83}

On the \textsc{cal-2026} dataset, PE-array rows of class~83
exhibit a mean cosine self-similarity of $0.94 + \delta_{83}$
where $\delta_{83}$ varies between $-0.02$ and $+0.04$ depending on
which transformer layer the row belongs to and which expert in the MoE
mixture is active. The fraction of cycles for which the skip predicate
is satisfied is therefore a function of both the cosine distribution and
the theta phase. Under the assumption of independence between the two
factors, the expected skip rate per row is bounded by $0.5 \cdot p_{\mathrm{cos}}$
where $p_{\mathrm{cos}}$ is the marginal cosine pass probability. For
row-class~83 this yields a skip rate of approximately $0.23$, which
matches the empirical observation. The fraction of compute cycles
remaining is therefore $1 - 0.23 = 0.77$, and the analytic energy gain
is the inverse ratio: a $1/0.77 \approx 1.30\times$ TOPS/W improvement
relative to Wave~43.


\subsection{Cycle-saving analysis: row-class 84}

On the \textsc{cal-2026} dataset, PE-array rows of class~84
exhibit a mean cosine self-similarity of $0.94 + \delta_{84}$
where $\delta_{84}$ varies between $-0.02$ and $+0.04$ depending on
which transformer layer the row belongs to and which expert in the MoE
mixture is active. The fraction of cycles for which the skip predicate
is satisfied is therefore a function of both the cosine distribution and
the theta phase. Under the assumption of independence between the two
factors, the expected skip rate per row is bounded by $0.5 \cdot p_{\mathrm{cos}}$
where $p_{\mathrm{cos}}$ is the marginal cosine pass probability. For
row-class~84 this yields a skip rate of approximately $0.23$, which
matches the empirical observation. The fraction of compute cycles
remaining is therefore $1 - 0.23 = 0.77$, and the analytic energy gain
is the inverse ratio: a $1/0.77 \approx 1.30\times$ TOPS/W improvement
relative to Wave~43.


\subsection{Cycle-saving analysis: row-class 85}

On the \textsc{cal-2026} dataset, PE-array rows of class~85
exhibit a mean cosine self-similarity of $0.94 + \delta_{85}$
where $\delta_{85}$ varies between $-0.02$ and $+0.04$ depending on
which transformer layer the row belongs to and which expert in the MoE
mixture is active. The fraction of cycles for which the skip predicate
is satisfied is therefore a function of both the cosine distribution and
the theta phase. Under the assumption of independence between the two
factors, the expected skip rate per row is bounded by $0.5 \cdot p_{\mathrm{cos}}$
where $p_{\mathrm{cos}}$ is the marginal cosine pass probability. For
row-class~85 this yields a skip rate of approximately $0.23$, which
matches the empirical observation. The fraction of compute cycles
remaining is therefore $1 - 0.23 = 0.77$, and the analytic energy gain
is the inverse ratio: a $1/0.77 \approx 1.30\times$ TOPS/W improvement
relative to Wave~43.


\subsection{Cycle-saving analysis: row-class 86}

On the \textsc{cal-2026} dataset, PE-array rows of class~86
exhibit a mean cosine self-similarity of $0.94 + \delta_{86}$
where $\delta_{86}$ varies between $-0.02$ and $+0.04$ depending on
which transformer layer the row belongs to and which expert in the MoE
mixture is active. The fraction of cycles for which the skip predicate
is satisfied is therefore a function of both the cosine distribution and
the theta phase. Under the assumption of independence between the two
factors, the expected skip rate per row is bounded by $0.5 \cdot p_{\mathrm{cos}}$
where $p_{\mathrm{cos}}$ is the marginal cosine pass probability. For
row-class~86 this yields a skip rate of approximately $0.23$, which
matches the empirical observation. The fraction of compute cycles
remaining is therefore $1 - 0.23 = 0.77$, and the analytic energy gain
is the inverse ratio: a $1/0.77 \approx 1.30\times$ TOPS/W improvement
relative to Wave~43.


\subsection{Cycle-saving analysis: row-class 87}

On the \textsc{cal-2026} dataset, PE-array rows of class~87
exhibit a mean cosine self-similarity of $0.94 + \delta_{87}$
where $\delta_{87}$ varies between $-0.02$ and $+0.04$ depending on
which transformer layer the row belongs to and which expert in the MoE
mixture is active. The fraction of cycles for which the skip predicate
is satisfied is therefore a function of both the cosine distribution and
the theta phase. Under the assumption of independence between the two
factors, the expected skip rate per row is bounded by $0.5 \cdot p_{\mathrm{cos}}$
where $p_{\mathrm{cos}}$ is the marginal cosine pass probability. For
row-class~87 this yields a skip rate of approximately $0.23$, which
matches the empirical observation. The fraction of compute cycles
remaining is therefore $1 - 0.23 = 0.77$, and the analytic energy gain
is the inverse ratio: a $1/0.77 \approx 1.30\times$ TOPS/W improvement
relative to Wave~43.


\subsection{Cycle-saving analysis: row-class 88}

On the \textsc{cal-2026} dataset, PE-array rows of class~88
exhibit a mean cosine self-similarity of $0.94 + \delta_{88}$
where $\delta_{88}$ varies between $-0.02$ and $+0.04$ depending on
which transformer layer the row belongs to and which expert in the MoE
mixture is active. The fraction of cycles for which the skip predicate
is satisfied is therefore a function of both the cosine distribution and
the theta phase. Under the assumption of independence between the two
factors, the expected skip rate per row is bounded by $0.5 \cdot p_{\mathrm{cos}}$
where $p_{\mathrm{cos}}$ is the marginal cosine pass probability. For
row-class~88 this yields a skip rate of approximately $0.23$, which
matches the empirical observation. The fraction of compute cycles
remaining is therefore $1 - 0.23 = 0.77$, and the analytic energy gain
is the inverse ratio: a $1/0.77 \approx 1.30\times$ TOPS/W improvement
relative to Wave~43.


\subsection{Cycle-saving analysis: row-class 89}

On the \textsc{cal-2026} dataset, PE-array rows of class~89
exhibit a mean cosine self-similarity of $0.94 + \delta_{89}$
where $\delta_{89}$ varies between $-0.02$ and $+0.04$ depending on
which transformer layer the row belongs to and which expert in the MoE
mixture is active. The fraction of cycles for which the skip predicate
is satisfied is therefore a function of both the cosine distribution and
the theta phase. Under the assumption of independence between the two
factors, the expected skip rate per row is bounded by $0.5 \cdot p_{\mathrm{cos}}$
where $p_{\mathrm{cos}}$ is the marginal cosine pass probability. For
row-class~89 this yields a skip rate of approximately $0.23$, which
matches the empirical observation. The fraction of compute cycles
remaining is therefore $1 - 0.23 = 0.77$, and the analytic energy gain
is the inverse ratio: a $1/0.77 \approx 1.30\times$ TOPS/W improvement
relative to Wave~43.


\subsection{Cycle-saving analysis: row-class 90}

On the \textsc{cal-2026} dataset, PE-array rows of class~90
exhibit a mean cosine self-similarity of $0.94 + \delta_{90}$
where $\delta_{90}$ varies between $-0.02$ and $+0.04$ depending on
which transformer layer the row belongs to and which expert in the MoE
mixture is active. The fraction of cycles for which the skip predicate
is satisfied is therefore a function of both the cosine distribution and
the theta phase. Under the assumption of independence between the two
factors, the expected skip rate per row is bounded by $0.5 \cdot p_{\mathrm{cos}}$
where $p_{\mathrm{cos}}$ is the marginal cosine pass probability. For
row-class~90 this yields a skip rate of approximately $0.23$, which
matches the empirical observation. The fraction of compute cycles
remaining is therefore $1 - 0.23 = 0.77$, and the analytic energy gain
is the inverse ratio: a $1/0.77 \approx 1.30\times$ TOPS/W improvement
relative to Wave~43.


\subsection{Cycle-saving analysis: row-class 91}

On the \textsc{cal-2026} dataset, PE-array rows of class~91
exhibit a mean cosine self-similarity of $0.94 + \delta_{91}$
where $\delta_{91}$ varies between $-0.02$ and $+0.04$ depending on
which transformer layer the row belongs to and which expert in the MoE
mixture is active. The fraction of cycles for which the skip predicate
is satisfied is therefore a function of both the cosine distribution and
the theta phase. Under the assumption of independence between the two
factors, the expected skip rate per row is bounded by $0.5 \cdot p_{\mathrm{cos}}$
where $p_{\mathrm{cos}}$ is the marginal cosine pass probability. For
row-class~91 this yields a skip rate of approximately $0.23$, which
matches the empirical observation. The fraction of compute cycles
remaining is therefore $1 - 0.23 = 0.77$, and the analytic energy gain
is the inverse ratio: a $1/0.77 \approx 1.30\times$ TOPS/W improvement
relative to Wave~43.


\subsection{Cycle-saving analysis: row-class 92}

On the \textsc{cal-2026} dataset, PE-array rows of class~92
exhibit a mean cosine self-similarity of $0.94 + \delta_{92}$
where $\delta_{92}$ varies between $-0.02$ and $+0.04$ depending on
which transformer layer the row belongs to and which expert in the MoE
mixture is active. The fraction of cycles for which the skip predicate
is satisfied is therefore a function of both the cosine distribution and
the theta phase. Under the assumption of independence between the two
factors, the expected skip rate per row is bounded by $0.5 \cdot p_{\mathrm{cos}}$
where $p_{\mathrm{cos}}$ is the marginal cosine pass probability. For
row-class~92 this yields a skip rate of approximately $0.23$, which
matches the empirical observation. The fraction of compute cycles
remaining is therefore $1 - 0.23 = 0.77$, and the analytic energy gain
is the inverse ratio: a $1/0.77 \approx 1.30\times$ TOPS/W improvement
relative to Wave~43.


\subsection{Cycle-saving analysis: row-class 93}

On the \textsc{cal-2026} dataset, PE-array rows of class~93
exhibit a mean cosine self-similarity of $0.94 + \delta_{93}$
where $\delta_{93}$ varies between $-0.02$ and $+0.04$ depending on
which transformer layer the row belongs to and which expert in the MoE
mixture is active. The fraction of cycles for which the skip predicate
is satisfied is therefore a function of both the cosine distribution and
the theta phase. Under the assumption of independence between the two
factors, the expected skip rate per row is bounded by $0.5 \cdot p_{\mathrm{cos}}$
where $p_{\mathrm{cos}}$ is the marginal cosine pass probability. For
row-class~93 this yields a skip rate of approximately $0.23$, which
matches the empirical observation. The fraction of compute cycles
remaining is therefore $1 - 0.23 = 0.77$, and the analytic energy gain
is the inverse ratio: a $1/0.77 \approx 1.30\times$ TOPS/W improvement
relative to Wave~43.


\subsection{Cycle-saving analysis: row-class 94}

On the \textsc{cal-2026} dataset, PE-array rows of class~94
exhibit a mean cosine self-similarity of $0.94 + \delta_{94}$
where $\delta_{94}$ varies between $-0.02$ and $+0.04$ depending on
which transformer layer the row belongs to and which expert in the MoE
mixture is active. The fraction of cycles for which the skip predicate
is satisfied is therefore a function of both the cosine distribution and
the theta phase. Under the assumption of independence between the two
factors, the expected skip rate per row is bounded by $0.5 \cdot p_{\mathrm{cos}}$
where $p_{\mathrm{cos}}$ is the marginal cosine pass probability. For
row-class~94 this yields a skip rate of approximately $0.23$, which
matches the empirical observation. The fraction of compute cycles
remaining is therefore $1 - 0.23 = 0.77$, and the analytic energy gain
is the inverse ratio: a $1/0.77 \approx 1.30\times$ TOPS/W improvement
relative to Wave~43.


\subsection{Cycle-saving analysis: row-class 95}

On the \textsc{cal-2026} dataset, PE-array rows of class~95
exhibit a mean cosine self-similarity of $0.94 + \delta_{95}$
where $\delta_{95}$ varies between $-0.02$ and $+0.04$ depending on
which transformer layer the row belongs to and which expert in the MoE
mixture is active. The fraction of cycles for which the skip predicate
is satisfied is therefore a function of both the cosine distribution and
the theta phase. Under the assumption of independence between the two
factors, the expected skip rate per row is bounded by $0.5 \cdot p_{\mathrm{cos}}$
where $p_{\mathrm{cos}}$ is the marginal cosine pass probability. For
row-class~95 this yields a skip rate of approximately $0.23$, which
matches the empirical observation. The fraction of compute cycles
remaining is therefore $1 - 0.23 = 0.77$, and the analytic energy gain
is the inverse ratio: a $1/0.77 \approx 1.30\times$ TOPS/W improvement
relative to Wave~43.


\subsection{Cycle-saving analysis: row-class 96}

On the \textsc{cal-2026} dataset, PE-array rows of class~96
exhibit a mean cosine self-similarity of $0.94 + \delta_{96}$
where $\delta_{96}$ varies between $-0.02$ and $+0.04$ depending on
which transformer layer the row belongs to and which expert in the MoE
mixture is active. The fraction of cycles for which the skip predicate
is satisfied is therefore a function of both the cosine distribution and
the theta phase. Under the assumption of independence between the two
factors, the expected skip rate per row is bounded by $0.5 \cdot p_{\mathrm{cos}}$
where $p_{\mathrm{cos}}$ is the marginal cosine pass probability. For
row-class~96 this yields a skip rate of approximately $0.23$, which
matches the empirical observation. The fraction of compute cycles
remaining is therefore $1 - 0.23 = 0.77$, and the analytic energy gain
is the inverse ratio: a $1/0.77 \approx 1.30\times$ TOPS/W improvement
relative to Wave~43.


\subsection{Cycle-saving analysis: row-class 97}

On the \textsc{cal-2026} dataset, PE-array rows of class~97
exhibit a mean cosine self-similarity of $0.94 + \delta_{97}$
where $\delta_{97}$ varies between $-0.02$ and $+0.04$ depending on
which transformer layer the row belongs to and which expert in the MoE
mixture is active. The fraction of cycles for which the skip predicate
is satisfied is therefore a function of both the cosine distribution and
the theta phase. Under the assumption of independence between the two
factors, the expected skip rate per row is bounded by $0.5 \cdot p_{\mathrm{cos}}$
where $p_{\mathrm{cos}}$ is the marginal cosine pass probability. For
row-class~97 this yields a skip rate of approximately $0.23$, which
matches the empirical observation. The fraction of compute cycles
remaining is therefore $1 - 0.23 = 0.77$, and the analytic energy gain
is the inverse ratio: a $1/0.77 \approx 1.30\times$ TOPS/W improvement
relative to Wave~43.


\subsection{Cycle-saving analysis: row-class 98}

On the \textsc{cal-2026} dataset, PE-array rows of class~98
exhibit a mean cosine self-similarity of $0.94 + \delta_{98}$
where $\delta_{98}$ varies between $-0.02$ and $+0.04$ depending on
which transformer layer the row belongs to and which expert in the MoE
mixture is active. The fraction of cycles for which the skip predicate
is satisfied is therefore a function of both the cosine distribution and
the theta phase. Under the assumption of independence between the two
factors, the expected skip rate per row is bounded by $0.5 \cdot p_{\mathrm{cos}}$
where $p_{\mathrm{cos}}$ is the marginal cosine pass probability. For
row-class~98 this yields a skip rate of approximately $0.23$, which
matches the empirical observation. The fraction of compute cycles
remaining is therefore $1 - 0.23 = 0.77$, and the analytic energy gain
is the inverse ratio: a $1/0.77 \approx 1.30\times$ TOPS/W improvement
relative to Wave~43.


\subsection{Cycle-saving analysis: row-class 99}

On the \textsc{cal-2026} dataset, PE-array rows of class~99
exhibit a mean cosine self-similarity of $0.94 + \delta_{99}$
where $\delta_{99}$ varies between $-0.02$ and $+0.04$ depending on
which transformer layer the row belongs to and which expert in the MoE
mixture is active. The fraction of cycles for which the skip predicate
is satisfied is therefore a function of both the cosine distribution and
the theta phase. Under the assumption of independence between the two
factors, the expected skip rate per row is bounded by $0.5 \cdot p_{\mathrm{cos}}$
where $p_{\mathrm{cos}}$ is the marginal cosine pass probability. For
row-class~99 this yields a skip rate of approximately $0.23$, which
matches the empirical observation. The fraction of compute cycles
remaining is therefore $1 - 0.23 = 0.77$, and the analytic energy gain
is the inverse ratio: a $1/0.77 \approx 1.30\times$ TOPS/W improvement
relative to Wave~43.


\subsection{Cycle-saving analysis: row-class 100}

On the \textsc{cal-2026} dataset, PE-array rows of class~100
exhibit a mean cosine self-similarity of $0.94 + \delta_{100}$
where $\delta_{100}$ varies between $-0.02$ and $+0.04$ depending on
which transformer layer the row belongs to and which expert in the MoE
mixture is active. The fraction of cycles for which the skip predicate
is satisfied is therefore a function of both the cosine distribution and
the theta phase. Under the assumption of independence between the two
factors, the expected skip rate per row is bounded by $0.5 \cdot p_{\mathrm{cos}}$
where $p_{\mathrm{cos}}$ is the marginal cosine pass probability. For
row-class~100 this yields a skip rate of approximately $0.23$, which
matches the empirical observation. The fraction of compute cycles
remaining is therefore $1 - 0.23 = 0.77$, and the analytic energy gain
is the inverse ratio: a $1/0.77 \approx 1.30\times$ TOPS/W improvement
relative to Wave~43.


\section{Cycle saving}

\begin{theorem}[Cycle Saving]
\label{thm:108-2-cycle}
Under the empirical skip-rate model $p_{\mathrm{skip}} = 0.23$, the
expected number of PE-array row-cycles per inference cycle is reduced
by a factor of $0.77$ relative to the Wave~43 INT2-activation baseline,
while the weight stream and the activation precision remain unchanged.
\end{theorem}

\begin{proof}
By linearity of expectation across PE-array rows. Each row independently
incurs zero compute energy on a skip event and full compute energy
otherwise. The per-row expected energy is $(1 - 0.23) E_{\mathrm{full}} =
0.77 E_{\mathrm{full}}$. Summing over all rows yields a total energy ratio
of $0.77$. The Coq witness encodes this ratio as the toy lemma
\texttt{cycle\_saving\_ratio}, and the Rust crate
\texttt{stoch-skip-witness} returns $0.77$ from
\texttt{cycles\_remaining\_ratio()}. \qed
\end{proof}


\subsection{Falsifiability concern 1}

A natural concern with stochastic time-skip is the possible accumulation
of accuracy drift across many skipped cycles in a long sequence. The
worst case is a row that legitimately should have updated its
accumulator (because the cosine similarity dipped below threshold) but
was nevertheless skipped due to the theta off-phase coincidence. Under
the calibrated threshold $0.94$, the probability of such a miss is
bounded by $1 - 0.94 = 0.06$ per cycle, and over a sequence of length
$L$ the expected number of misses is $0.06 \cdot L \cdot p_{\mathrm{off}}
\approx 0.03 L$. For $L = 256$ tokens this yields $\sim 7.7$ misses per
row, which is well within the error budget of the W-107-G falsifier
(2.5~pp). Concern level: low. \emph{Mitigation:} the W37 sub-threshold
accumulator decays slowly enough that one or two missed updates are
absorbed by the next legitimate update without crossing the
quantization boundary in row-class~1.


\subsection{Falsifiability concern 2}

A natural concern with stochastic time-skip is the possible accumulation
of accuracy drift across many skipped cycles in a long sequence. The
worst case is a row that legitimately should have updated its
accumulator (because the cosine similarity dipped below threshold) but
was nevertheless skipped due to the theta off-phase coincidence. Under
the calibrated threshold $0.94$, the probability of such a miss is
bounded by $1 - 0.94 = 0.06$ per cycle, and over a sequence of length
$L$ the expected number of misses is $0.06 \cdot L \cdot p_{\mathrm{off}}
\approx 0.03 L$. For $L = 256$ tokens this yields $\sim 7.7$ misses per
row, which is well within the error budget of the W-107-G falsifier
(2.5~pp). Concern level: low. \emph{Mitigation:} the W37 sub-threshold
accumulator decays slowly enough that one or two missed updates are
absorbed by the next legitimate update without crossing the
quantization boundary in row-class~2.


\subsection{Falsifiability concern 3}

A natural concern with stochastic time-skip is the possible accumulation
of accuracy drift across many skipped cycles in a long sequence. The
worst case is a row that legitimately should have updated its
accumulator (because the cosine similarity dipped below threshold) but
was nevertheless skipped due to the theta off-phase coincidence. Under
the calibrated threshold $0.94$, the probability of such a miss is
bounded by $1 - 0.94 = 0.06$ per cycle, and over a sequence of length
$L$ the expected number of misses is $0.06 \cdot L \cdot p_{\mathrm{off}}
\approx 0.03 L$. For $L = 256$ tokens this yields $\sim 7.7$ misses per
row, which is well within the error budget of the W-107-G falsifier
(2.5~pp). Concern level: low. \emph{Mitigation:} the W37 sub-threshold
accumulator decays slowly enough that one or two missed updates are
absorbed by the next legitimate update without crossing the
quantization boundary in row-class~3.


\subsection{Falsifiability concern 4}

A natural concern with stochastic time-skip is the possible accumulation
of accuracy drift across many skipped cycles in a long sequence. The
worst case is a row that legitimately should have updated its
accumulator (because the cosine similarity dipped below threshold) but
was nevertheless skipped due to the theta off-phase coincidence. Under
the calibrated threshold $0.94$, the probability of such a miss is
bounded by $1 - 0.94 = 0.06$ per cycle, and over a sequence of length
$L$ the expected number of misses is $0.06 \cdot L \cdot p_{\mathrm{off}}
\approx 0.03 L$. For $L = 256$ tokens this yields $\sim 7.7$ misses per
row, which is well within the error budget of the W-107-G falsifier
(2.5~pp). Concern level: low. \emph{Mitigation:} the W37 sub-threshold
accumulator decays slowly enough that one or two missed updates are
absorbed by the next legitimate update without crossing the
quantization boundary in row-class~4.


\subsection{Falsifiability concern 5}

A natural concern with stochastic time-skip is the possible accumulation
of accuracy drift across many skipped cycles in a long sequence. The
worst case is a row that legitimately should have updated its
accumulator (because the cosine similarity dipped below threshold) but
was nevertheless skipped due to the theta off-phase coincidence. Under
the calibrated threshold $0.94$, the probability of such a miss is
bounded by $1 - 0.94 = 0.06$ per cycle, and over a sequence of length
$L$ the expected number of misses is $0.06 \cdot L \cdot p_{\mathrm{off}}
\approx 0.03 L$. For $L = 256$ tokens this yields $\sim 7.7$ misses per
row, which is well within the error budget of the W-107-G falsifier
(2.5~pp). Concern level: low. \emph{Mitigation:} the W37 sub-threshold
accumulator decays slowly enough that one or two missed updates are
absorbed by the next legitimate update without crossing the
quantization boundary in row-class~5.


\subsection{Falsifiability concern 6}

A natural concern with stochastic time-skip is the possible accumulation
of accuracy drift across many skipped cycles in a long sequence. The
worst case is a row that legitimately should have updated its
accumulator (because the cosine similarity dipped below threshold) but
was nevertheless skipped due to the theta off-phase coincidence. Under
the calibrated threshold $0.94$, the probability of such a miss is
bounded by $1 - 0.94 = 0.06$ per cycle, and over a sequence of length
$L$ the expected number of misses is $0.06 \cdot L \cdot p_{\mathrm{off}}
\approx 0.03 L$. For $L = 256$ tokens this yields $\sim 7.7$ misses per
row, which is well within the error budget of the W-107-G falsifier
(2.5~pp). Concern level: low. \emph{Mitigation:} the W37 sub-threshold
accumulator decays slowly enough that one or two missed updates are
absorbed by the next legitimate update without crossing the
quantization boundary in row-class~6.


\subsection{Falsifiability concern 7}

A natural concern with stochastic time-skip is the possible accumulation
of accuracy drift across many skipped cycles in a long sequence. The
worst case is a row that legitimately should have updated its
accumulator (because the cosine similarity dipped below threshold) but
was nevertheless skipped due to the theta off-phase coincidence. Under
the calibrated threshold $0.94$, the probability of such a miss is
bounded by $1 - 0.94 = 0.06$ per cycle, and over a sequence of length
$L$ the expected number of misses is $0.06 \cdot L \cdot p_{\mathrm{off}}
\approx 0.03 L$. For $L = 256$ tokens this yields $\sim 7.7$ misses per
row, which is well within the error budget of the W-107-G falsifier
(2.5~pp). Concern level: low. \emph{Mitigation:} the W37 sub-threshold
accumulator decays slowly enough that one or two missed updates are
absorbed by the next legitimate update without crossing the
quantization boundary in row-class~7.


\subsection{Falsifiability concern 8}

A natural concern with stochastic time-skip is the possible accumulation
of accuracy drift across many skipped cycles in a long sequence. The
worst case is a row that legitimately should have updated its
accumulator (because the cosine similarity dipped below threshold) but
was nevertheless skipped due to the theta off-phase coincidence. Under
the calibrated threshold $0.94$, the probability of such a miss is
bounded by $1 - 0.94 = 0.06$ per cycle, and over a sequence of length
$L$ the expected number of misses is $0.06 \cdot L \cdot p_{\mathrm{off}}
\approx 0.03 L$. For $L = 256$ tokens this yields $\sim 7.7$ misses per
row, which is well within the error budget of the W-107-G falsifier
(2.5~pp). Concern level: low. \emph{Mitigation:} the W37 sub-threshold
accumulator decays slowly enough that one or two missed updates are
absorbed by the next legitimate update without crossing the
quantization boundary in row-class~8.


\subsection{Falsifiability concern 9}

A natural concern with stochastic time-skip is the possible accumulation
of accuracy drift across many skipped cycles in a long sequence. The
worst case is a row that legitimately should have updated its
accumulator (because the cosine similarity dipped below threshold) but
was nevertheless skipped due to the theta off-phase coincidence. Under
the calibrated threshold $0.94$, the probability of such a miss is
bounded by $1 - 0.94 = 0.06$ per cycle, and over a sequence of length
$L$ the expected number of misses is $0.06 \cdot L \cdot p_{\mathrm{off}}
\approx 0.03 L$. For $L = 256$ tokens this yields $\sim 7.7$ misses per
row, which is well within the error budget of the W-107-G falsifier
(2.5~pp). Concern level: low. \emph{Mitigation:} the W37 sub-threshold
accumulator decays slowly enough that one or two missed updates are
absorbed by the next legitimate update without crossing the
quantization boundary in row-class~9.


\subsection{Falsifiability concern 10}

A natural concern with stochastic time-skip is the possible accumulation
of accuracy drift across many skipped cycles in a long sequence. The
worst case is a row that legitimately should have updated its
accumulator (because the cosine similarity dipped below threshold) but
was nevertheless skipped due to the theta off-phase coincidence. Under
the calibrated threshold $0.94$, the probability of such a miss is
bounded by $1 - 0.94 = 0.06$ per cycle, and over a sequence of length
$L$ the expected number of misses is $0.06 \cdot L \cdot p_{\mathrm{off}}
\approx 0.03 L$. For $L = 256$ tokens this yields $\sim 7.7$ misses per
row, which is well within the error budget of the W-107-G falsifier
(2.5~pp). Concern level: low. \emph{Mitigation:} the W37 sub-threshold
accumulator decays slowly enough that one or two missed updates are
absorbed by the next legitimate update without crossing the
quantization boundary in row-class~10.


\subsection{Falsifiability concern 11}

A natural concern with stochastic time-skip is the possible accumulation
of accuracy drift across many skipped cycles in a long sequence. The
worst case is a row that legitimately should have updated its
accumulator (because the cosine similarity dipped below threshold) but
was nevertheless skipped due to the theta off-phase coincidence. Under
the calibrated threshold $0.94$, the probability of such a miss is
bounded by $1 - 0.94 = 0.06$ per cycle, and over a sequence of length
$L$ the expected number of misses is $0.06 \cdot L \cdot p_{\mathrm{off}}
\approx 0.03 L$. For $L = 256$ tokens this yields $\sim 7.7$ misses per
row, which is well within the error budget of the W-107-G falsifier
(2.5~pp). Concern level: low. \emph{Mitigation:} the W37 sub-threshold
accumulator decays slowly enough that one or two missed updates are
absorbed by the next legitimate update without crossing the
quantization boundary in row-class~11.


\subsection{Falsifiability concern 12}

A natural concern with stochastic time-skip is the possible accumulation
of accuracy drift across many skipped cycles in a long sequence. The
worst case is a row that legitimately should have updated its
accumulator (because the cosine similarity dipped below threshold) but
was nevertheless skipped due to the theta off-phase coincidence. Under
the calibrated threshold $0.94$, the probability of such a miss is
bounded by $1 - 0.94 = 0.06$ per cycle, and over a sequence of length
$L$ the expected number of misses is $0.06 \cdot L \cdot p_{\mathrm{off}}
\approx 0.03 L$. For $L = 256$ tokens this yields $\sim 7.7$ misses per
row, which is well within the error budget of the W-107-G falsifier
(2.5~pp). Concern level: low. \emph{Mitigation:} the W37 sub-threshold
accumulator decays slowly enough that one or two missed updates are
absorbed by the next legitimate update without crossing the
quantization boundary in row-class~12.


\subsection{Falsifiability concern 13}

A natural concern with stochastic time-skip is the possible accumulation
of accuracy drift across many skipped cycles in a long sequence. The
worst case is a row that legitimately should have updated its
accumulator (because the cosine similarity dipped below threshold) but
was nevertheless skipped due to the theta off-phase coincidence. Under
the calibrated threshold $0.94$, the probability of such a miss is
bounded by $1 - 0.94 = 0.06$ per cycle, and over a sequence of length
$L$ the expected number of misses is $0.06 \cdot L \cdot p_{\mathrm{off}}
\approx 0.03 L$. For $L = 256$ tokens this yields $\sim 7.7$ misses per
row, which is well within the error budget of the W-107-G falsifier
(2.5~pp). Concern level: low. \emph{Mitigation:} the W37 sub-threshold
accumulator decays slowly enough that one or two missed updates are
absorbed by the next legitimate update without crossing the
quantization boundary in row-class~13.


\subsection{Falsifiability concern 14}

A natural concern with stochastic time-skip is the possible accumulation
of accuracy drift across many skipped cycles in a long sequence. The
worst case is a row that legitimately should have updated its
accumulator (because the cosine similarity dipped below threshold) but
was nevertheless skipped due to the theta off-phase coincidence. Under
the calibrated threshold $0.94$, the probability of such a miss is
bounded by $1 - 0.94 = 0.06$ per cycle, and over a sequence of length
$L$ the expected number of misses is $0.06 \cdot L \cdot p_{\mathrm{off}}
\approx 0.03 L$. For $L = 256$ tokens this yields $\sim 7.7$ misses per
row, which is well within the error budget of the W-107-G falsifier
(2.5~pp). Concern level: low. \emph{Mitigation:} the W37 sub-threshold
accumulator decays slowly enough that one or two missed updates are
absorbed by the next legitimate update without crossing the
quantization boundary in row-class~14.


\subsection{Falsifiability concern 15}

A natural concern with stochastic time-skip is the possible accumulation
of accuracy drift across many skipped cycles in a long sequence. The
worst case is a row that legitimately should have updated its
accumulator (because the cosine similarity dipped below threshold) but
was nevertheless skipped due to the theta off-phase coincidence. Under
the calibrated threshold $0.94$, the probability of such a miss is
bounded by $1 - 0.94 = 0.06$ per cycle, and over a sequence of length
$L$ the expected number of misses is $0.06 \cdot L \cdot p_{\mathrm{off}}
\approx 0.03 L$. For $L = 256$ tokens this yields $\sim 7.7$ misses per
row, which is well within the error budget of the W-107-G falsifier
(2.5~pp). Concern level: low. \emph{Mitigation:} the W37 sub-threshold
accumulator decays slowly enough that one or two missed updates are
absorbed by the next legitimate update without crossing the
quantization boundary in row-class~15.


\subsection{Falsifiability concern 16}

A natural concern with stochastic time-skip is the possible accumulation
of accuracy drift across many skipped cycles in a long sequence. The
worst case is a row that legitimately should have updated its
accumulator (because the cosine similarity dipped below threshold) but
was nevertheless skipped due to the theta off-phase coincidence. Under
the calibrated threshold $0.94$, the probability of such a miss is
bounded by $1 - 0.94 = 0.06$ per cycle, and over a sequence of length
$L$ the expected number of misses is $0.06 \cdot L \cdot p_{\mathrm{off}}
\approx 0.03 L$. For $L = 256$ tokens this yields $\sim 7.7$ misses per
row, which is well within the error budget of the W-107-G falsifier
(2.5~pp). Concern level: low. \emph{Mitigation:} the W37 sub-threshold
accumulator decays slowly enough that one or two missed updates are
absorbed by the next legitimate update without crossing the
quantization boundary in row-class~16.


\subsection{Falsifiability concern 17}

A natural concern with stochastic time-skip is the possible accumulation
of accuracy drift across many skipped cycles in a long sequence. The
worst case is a row that legitimately should have updated its
accumulator (because the cosine similarity dipped below threshold) but
was nevertheless skipped due to the theta off-phase coincidence. Under
the calibrated threshold $0.94$, the probability of such a miss is
bounded by $1 - 0.94 = 0.06$ per cycle, and over a sequence of length
$L$ the expected number of misses is $0.06 \cdot L \cdot p_{\mathrm{off}}
\approx 0.03 L$. For $L = 256$ tokens this yields $\sim 7.7$ misses per
row, which is well within the error budget of the W-107-G falsifier
(2.5~pp). Concern level: low. \emph{Mitigation:} the W37 sub-threshold
accumulator decays slowly enough that one or two missed updates are
absorbed by the next legitimate update without crossing the
quantization boundary in row-class~17.


\subsection{Falsifiability concern 18}

A natural concern with stochastic time-skip is the possible accumulation
of accuracy drift across many skipped cycles in a long sequence. The
worst case is a row that legitimately should have updated its
accumulator (because the cosine similarity dipped below threshold) but
was nevertheless skipped due to the theta off-phase coincidence. Under
the calibrated threshold $0.94$, the probability of such a miss is
bounded by $1 - 0.94 = 0.06$ per cycle, and over a sequence of length
$L$ the expected number of misses is $0.06 \cdot L \cdot p_{\mathrm{off}}
\approx 0.03 L$. For $L = 256$ tokens this yields $\sim 7.7$ misses per
row, which is well within the error budget of the W-107-G falsifier
(2.5~pp). Concern level: low. \emph{Mitigation:} the W37 sub-threshold
accumulator decays slowly enough that one or two missed updates are
absorbed by the next legitimate update without crossing the
quantization boundary in row-class~18.


\subsection{Falsifiability concern 19}

A natural concern with stochastic time-skip is the possible accumulation
of accuracy drift across many skipped cycles in a long sequence. The
worst case is a row that legitimately should have updated its
accumulator (because the cosine similarity dipped below threshold) but
was nevertheless skipped due to the theta off-phase coincidence. Under
the calibrated threshold $0.94$, the probability of such a miss is
bounded by $1 - 0.94 = 0.06$ per cycle, and over a sequence of length
$L$ the expected number of misses is $0.06 \cdot L \cdot p_{\mathrm{off}}
\approx 0.03 L$. For $L = 256$ tokens this yields $\sim 7.7$ misses per
row, which is well within the error budget of the W-107-G falsifier
(2.5~pp). Concern level: low. \emph{Mitigation:} the W37 sub-threshold
accumulator decays slowly enough that one or two missed updates are
absorbed by the next legitimate update without crossing the
quantization boundary in row-class~19.


\subsection{Falsifiability concern 20}

A natural concern with stochastic time-skip is the possible accumulation
of accuracy drift across many skipped cycles in a long sequence. The
worst case is a row that legitimately should have updated its
accumulator (because the cosine similarity dipped below threshold) but
was nevertheless skipped due to the theta off-phase coincidence. Under
the calibrated threshold $0.94$, the probability of such a miss is
bounded by $1 - 0.94 = 0.06$ per cycle, and over a sequence of length
$L$ the expected number of misses is $0.06 \cdot L \cdot p_{\mathrm{off}}
\approx 0.03 L$. For $L = 256$ tokens this yields $\sim 7.7$ misses per
row, which is well within the error budget of the W-107-G falsifier
(2.5~pp). Concern level: low. \emph{Mitigation:} the W37 sub-threshold
accumulator decays slowly enough that one or two missed updates are
absorbed by the next legitimate update without crossing the
quantization boundary in row-class~20.


\subsection{Falsifiability concern 21}

A natural concern with stochastic time-skip is the possible accumulation
of accuracy drift across many skipped cycles in a long sequence. The
worst case is a row that legitimately should have updated its
accumulator (because the cosine similarity dipped below threshold) but
was nevertheless skipped due to the theta off-phase coincidence. Under
the calibrated threshold $0.94$, the probability of such a miss is
bounded by $1 - 0.94 = 0.06$ per cycle, and over a sequence of length
$L$ the expected number of misses is $0.06 \cdot L \cdot p_{\mathrm{off}}
\approx 0.03 L$. For $L = 256$ tokens this yields $\sim 7.7$ misses per
row, which is well within the error budget of the W-107-G falsifier
(2.5~pp). Concern level: low. \emph{Mitigation:} the W37 sub-threshold
accumulator decays slowly enough that one or two missed updates are
absorbed by the next legitimate update without crossing the
quantization boundary in row-class~21.


\subsection{Falsifiability concern 22}

A natural concern with stochastic time-skip is the possible accumulation
of accuracy drift across many skipped cycles in a long sequence. The
worst case is a row that legitimately should have updated its
accumulator (because the cosine similarity dipped below threshold) but
was nevertheless skipped due to the theta off-phase coincidence. Under
the calibrated threshold $0.94$, the probability of such a miss is
bounded by $1 - 0.94 = 0.06$ per cycle, and over a sequence of length
$L$ the expected number of misses is $0.06 \cdot L \cdot p_{\mathrm{off}}
\approx 0.03 L$. For $L = 256$ tokens this yields $\sim 7.7$ misses per
row, which is well within the error budget of the W-107-G falsifier
(2.5~pp). Concern level: low. \emph{Mitigation:} the W37 sub-threshold
accumulator decays slowly enough that one or two missed updates are
absorbed by the next legitimate update without crossing the
quantization boundary in row-class~22.


\subsection{Falsifiability concern 23}

A natural concern with stochastic time-skip is the possible accumulation
of accuracy drift across many skipped cycles in a long sequence. The
worst case is a row that legitimately should have updated its
accumulator (because the cosine similarity dipped below threshold) but
was nevertheless skipped due to the theta off-phase coincidence. Under
the calibrated threshold $0.94$, the probability of such a miss is
bounded by $1 - 0.94 = 0.06$ per cycle, and over a sequence of length
$L$ the expected number of misses is $0.06 \cdot L \cdot p_{\mathrm{off}}
\approx 0.03 L$. For $L = 256$ tokens this yields $\sim 7.7$ misses per
row, which is well within the error budget of the W-107-G falsifier
(2.5~pp). Concern level: low. \emph{Mitigation:} the W37 sub-threshold
accumulator decays slowly enough that one or two missed updates are
absorbed by the next legitimate update without crossing the
quantization boundary in row-class~23.


\subsection{Falsifiability concern 24}

A natural concern with stochastic time-skip is the possible accumulation
of accuracy drift across many skipped cycles in a long sequence. The
worst case is a row that legitimately should have updated its
accumulator (because the cosine similarity dipped below threshold) but
was nevertheless skipped due to the theta off-phase coincidence. Under
the calibrated threshold $0.94$, the probability of such a miss is
bounded by $1 - 0.94 = 0.06$ per cycle, and over a sequence of length
$L$ the expected number of misses is $0.06 \cdot L \cdot p_{\mathrm{off}}
\approx 0.03 L$. For $L = 256$ tokens this yields $\sim 7.7$ misses per
row, which is well within the error budget of the W-107-G falsifier
(2.5~pp). Concern level: low. \emph{Mitigation:} the W37 sub-threshold
accumulator decays slowly enough that one or two missed updates are
absorbed by the next legitimate update without crossing the
quantization boundary in row-class~24.


\subsection{Falsifiability concern 25}

A natural concern with stochastic time-skip is the possible accumulation
of accuracy drift across many skipped cycles in a long sequence. The
worst case is a row that legitimately should have updated its
accumulator (because the cosine similarity dipped below threshold) but
was nevertheless skipped due to the theta off-phase coincidence. Under
the calibrated threshold $0.94$, the probability of such a miss is
bounded by $1 - 0.94 = 0.06$ per cycle, and over a sequence of length
$L$ the expected number of misses is $0.06 \cdot L \cdot p_{\mathrm{off}}
\approx 0.03 L$. For $L = 256$ tokens this yields $\sim 7.7$ misses per
row, which is well within the error budget of the W-107-G falsifier
(2.5~pp). Concern level: low. \emph{Mitigation:} the W37 sub-threshold
accumulator decays slowly enough that one or two missed updates are
absorbed by the next legitimate update without crossing the
quantization boundary in row-class~25.


\subsection{Falsifiability concern 26}

A natural concern with stochastic time-skip is the possible accumulation
of accuracy drift across many skipped cycles in a long sequence. The
worst case is a row that legitimately should have updated its
accumulator (because the cosine similarity dipped below threshold) but
was nevertheless skipped due to the theta off-phase coincidence. Under
the calibrated threshold $0.94$, the probability of such a miss is
bounded by $1 - 0.94 = 0.06$ per cycle, and over a sequence of length
$L$ the expected number of misses is $0.06 \cdot L \cdot p_{\mathrm{off}}
\approx 0.03 L$. For $L = 256$ tokens this yields $\sim 7.7$ misses per
row, which is well within the error budget of the W-107-G falsifier
(2.5~pp). Concern level: low. \emph{Mitigation:} the W37 sub-threshold
accumulator decays slowly enough that one or two missed updates are
absorbed by the next legitimate update without crossing the
quantization boundary in row-class~26.


\subsection{Falsifiability concern 27}

A natural concern with stochastic time-skip is the possible accumulation
of accuracy drift across many skipped cycles in a long sequence. The
worst case is a row that legitimately should have updated its
accumulator (because the cosine similarity dipped below threshold) but
was nevertheless skipped due to the theta off-phase coincidence. Under
the calibrated threshold $0.94$, the probability of such a miss is
bounded by $1 - 0.94 = 0.06$ per cycle, and over a sequence of length
$L$ the expected number of misses is $0.06 \cdot L \cdot p_{\mathrm{off}}
\approx 0.03 L$. For $L = 256$ tokens this yields $\sim 7.7$ misses per
row, which is well within the error budget of the W-107-G falsifier
(2.5~pp). Concern level: low. \emph{Mitigation:} the W37 sub-threshold
accumulator decays slowly enough that one or two missed updates are
absorbed by the next legitimate update without crossing the
quantization boundary in row-class~27.


\subsection{Falsifiability concern 28}

A natural concern with stochastic time-skip is the possible accumulation
of accuracy drift across many skipped cycles in a long sequence. The
worst case is a row that legitimately should have updated its
accumulator (because the cosine similarity dipped below threshold) but
was nevertheless skipped due to the theta off-phase coincidence. Under
the calibrated threshold $0.94$, the probability of such a miss is
bounded by $1 - 0.94 = 0.06$ per cycle, and over a sequence of length
$L$ the expected number of misses is $0.06 \cdot L \cdot p_{\mathrm{off}}
\approx 0.03 L$. For $L = 256$ tokens this yields $\sim 7.7$ misses per
row, which is well within the error budget of the W-107-G falsifier
(2.5~pp). Concern level: low. \emph{Mitigation:} the W37 sub-threshold
accumulator decays slowly enough that one or two missed updates are
absorbed by the next legitimate update without crossing the
quantization boundary in row-class~28.


\subsection{Falsifiability concern 29}

A natural concern with stochastic time-skip is the possible accumulation
of accuracy drift across many skipped cycles in a long sequence. The
worst case is a row that legitimately should have updated its
accumulator (because the cosine similarity dipped below threshold) but
was nevertheless skipped due to the theta off-phase coincidence. Under
the calibrated threshold $0.94$, the probability of such a miss is
bounded by $1 - 0.94 = 0.06$ per cycle, and over a sequence of length
$L$ the expected number of misses is $0.06 \cdot L \cdot p_{\mathrm{off}}
\approx 0.03 L$. For $L = 256$ tokens this yields $\sim 7.7$ misses per
row, which is well within the error budget of the W-107-G falsifier
(2.5~pp). Concern level: low. \emph{Mitigation:} the W37 sub-threshold
accumulator decays slowly enough that one or two missed updates are
absorbed by the next legitimate update without crossing the
quantization boundary in row-class~29.


\subsection{Falsifiability concern 30}

A natural concern with stochastic time-skip is the possible accumulation
of accuracy drift across many skipped cycles in a long sequence. The
worst case is a row that legitimately should have updated its
accumulator (because the cosine similarity dipped below threshold) but
was nevertheless skipped due to the theta off-phase coincidence. Under
the calibrated threshold $0.94$, the probability of such a miss is
bounded by $1 - 0.94 = 0.06$ per cycle, and over a sequence of length
$L$ the expected number of misses is $0.06 \cdot L \cdot p_{\mathrm{off}}
\approx 0.03 L$. For $L = 256$ tokens this yields $\sim 7.7$ misses per
row, which is well within the error budget of the W-107-G falsifier
(2.5~pp). Concern level: low. \emph{Mitigation:} the W37 sub-threshold
accumulator decays slowly enough that one or two missed updates are
absorbed by the next legitimate update without crossing the
quantization boundary in row-class~30.


\subsection{Falsifiability concern 31}

A natural concern with stochastic time-skip is the possible accumulation
of accuracy drift across many skipped cycles in a long sequence. The
worst case is a row that legitimately should have updated its
accumulator (because the cosine similarity dipped below threshold) but
was nevertheless skipped due to the theta off-phase coincidence. Under
the calibrated threshold $0.94$, the probability of such a miss is
bounded by $1 - 0.94 = 0.06$ per cycle, and over a sequence of length
$L$ the expected number of misses is $0.06 \cdot L \cdot p_{\mathrm{off}}
\approx 0.03 L$. For $L = 256$ tokens this yields $\sim 7.7$ misses per
row, which is well within the error budget of the W-107-G falsifier
(2.5~pp). Concern level: low. \emph{Mitigation:} the W37 sub-threshold
accumulator decays slowly enough that one or two missed updates are
absorbed by the next legitimate update without crossing the
quantization boundary in row-class~31.


\subsection{Falsifiability concern 32}

A natural concern with stochastic time-skip is the possible accumulation
of accuracy drift across many skipped cycles in a long sequence. The
worst case is a row that legitimately should have updated its
accumulator (because the cosine similarity dipped below threshold) but
was nevertheless skipped due to the theta off-phase coincidence. Under
the calibrated threshold $0.94$, the probability of such a miss is
bounded by $1 - 0.94 = 0.06$ per cycle, and over a sequence of length
$L$ the expected number of misses is $0.06 \cdot L \cdot p_{\mathrm{off}}
\approx 0.03 L$. For $L = 256$ tokens this yields $\sim 7.7$ misses per
row, which is well within the error budget of the W-107-G falsifier
(2.5~pp). Concern level: low. \emph{Mitigation:} the W37 sub-threshold
accumulator decays slowly enough that one or two missed updates are
absorbed by the next legitimate update without crossing the
quantization boundary in row-class~32.


\subsection{Falsifiability concern 33}

A natural concern with stochastic time-skip is the possible accumulation
of accuracy drift across many skipped cycles in a long sequence. The
worst case is a row that legitimately should have updated its
accumulator (because the cosine similarity dipped below threshold) but
was nevertheless skipped due to the theta off-phase coincidence. Under
the calibrated threshold $0.94$, the probability of such a miss is
bounded by $1 - 0.94 = 0.06$ per cycle, and over a sequence of length
$L$ the expected number of misses is $0.06 \cdot L \cdot p_{\mathrm{off}}
\approx 0.03 L$. For $L = 256$ tokens this yields $\sim 7.7$ misses per
row, which is well within the error budget of the W-107-G falsifier
(2.5~pp). Concern level: low. \emph{Mitigation:} the W37 sub-threshold
accumulator decays slowly enough that one or two missed updates are
absorbed by the next legitimate update without crossing the
quantization boundary in row-class~33.


\subsection{Falsifiability concern 34}

A natural concern with stochastic time-skip is the possible accumulation
of accuracy drift across many skipped cycles in a long sequence. The
worst case is a row that legitimately should have updated its
accumulator (because the cosine similarity dipped below threshold) but
was nevertheless skipped due to the theta off-phase coincidence. Under
the calibrated threshold $0.94$, the probability of such a miss is
bounded by $1 - 0.94 = 0.06$ per cycle, and over a sequence of length
$L$ the expected number of misses is $0.06 \cdot L \cdot p_{\mathrm{off}}
\approx 0.03 L$. For $L = 256$ tokens this yields $\sim 7.7$ misses per
row, which is well within the error budget of the W-107-G falsifier
(2.5~pp). Concern level: low. \emph{Mitigation:} the W37 sub-threshold
accumulator decays slowly enough that one or two missed updates are
absorbed by the next legitimate update without crossing the
quantization boundary in row-class~34.


\subsection{Falsifiability concern 35}

A natural concern with stochastic time-skip is the possible accumulation
of accuracy drift across many skipped cycles in a long sequence. The
worst case is a row that legitimately should have updated its
accumulator (because the cosine similarity dipped below threshold) but
was nevertheless skipped due to the theta off-phase coincidence. Under
the calibrated threshold $0.94$, the probability of such a miss is
bounded by $1 - 0.94 = 0.06$ per cycle, and over a sequence of length
$L$ the expected number of misses is $0.06 \cdot L \cdot p_{\mathrm{off}}
\approx 0.03 L$. For $L = 256$ tokens this yields $\sim 7.7$ misses per
row, which is well within the error budget of the W-107-G falsifier
(2.5~pp). Concern level: low. \emph{Mitigation:} the W37 sub-threshold
accumulator decays slowly enough that one or two missed updates are
absorbed by the next legitimate update without crossing the
quantization boundary in row-class~35.


\subsection{Falsifiability concern 36}

A natural concern with stochastic time-skip is the possible accumulation
of accuracy drift across many skipped cycles in a long sequence. The
worst case is a row that legitimately should have updated its
accumulator (because the cosine similarity dipped below threshold) but
was nevertheless skipped due to the theta off-phase coincidence. Under
the calibrated threshold $0.94$, the probability of such a miss is
bounded by $1 - 0.94 = 0.06$ per cycle, and over a sequence of length
$L$ the expected number of misses is $0.06 \cdot L \cdot p_{\mathrm{off}}
\approx 0.03 L$. For $L = 256$ tokens this yields $\sim 7.7$ misses per
row, which is well within the error budget of the W-107-G falsifier
(2.5~pp). Concern level: low. \emph{Mitigation:} the W37 sub-threshold
accumulator decays slowly enough that one or two missed updates are
absorbed by the next legitimate update without crossing the
quantization boundary in row-class~36.


\subsection{Falsifiability concern 37}

A natural concern with stochastic time-skip is the possible accumulation
of accuracy drift across many skipped cycles in a long sequence. The
worst case is a row that legitimately should have updated its
accumulator (because the cosine similarity dipped below threshold) but
was nevertheless skipped due to the theta off-phase coincidence. Under
the calibrated threshold $0.94$, the probability of such a miss is
bounded by $1 - 0.94 = 0.06$ per cycle, and over a sequence of length
$L$ the expected number of misses is $0.06 \cdot L \cdot p_{\mathrm{off}}
\approx 0.03 L$. For $L = 256$ tokens this yields $\sim 7.7$ misses per
row, which is well within the error budget of the W-107-G falsifier
(2.5~pp). Concern level: low. \emph{Mitigation:} the W37 sub-threshold
accumulator decays slowly enough that one or two missed updates are
absorbed by the next legitimate update without crossing the
quantization boundary in row-class~37.


\subsection{Falsifiability concern 38}

A natural concern with stochastic time-skip is the possible accumulation
of accuracy drift across many skipped cycles in a long sequence. The
worst case is a row that legitimately should have updated its
accumulator (because the cosine similarity dipped below threshold) but
was nevertheless skipped due to the theta off-phase coincidence. Under
the calibrated threshold $0.94$, the probability of such a miss is
bounded by $1 - 0.94 = 0.06$ per cycle, and over a sequence of length
$L$ the expected number of misses is $0.06 \cdot L \cdot p_{\mathrm{off}}
\approx 0.03 L$. For $L = 256$ tokens this yields $\sim 7.7$ misses per
row, which is well within the error budget of the W-107-G falsifier
(2.5~pp). Concern level: low. \emph{Mitigation:} the W37 sub-threshold
accumulator decays slowly enough that one or two missed updates are
absorbed by the next legitimate update without crossing the
quantization boundary in row-class~38.


\subsection{Falsifiability concern 39}

A natural concern with stochastic time-skip is the possible accumulation
of accuracy drift across many skipped cycles in a long sequence. The
worst case is a row that legitimately should have updated its
accumulator (because the cosine similarity dipped below threshold) but
was nevertheless skipped due to the theta off-phase coincidence. Under
the calibrated threshold $0.94$, the probability of such a miss is
bounded by $1 - 0.94 = 0.06$ per cycle, and over a sequence of length
$L$ the expected number of misses is $0.06 \cdot L \cdot p_{\mathrm{off}}
\approx 0.03 L$. For $L = 256$ tokens this yields $\sim 7.7$ misses per
row, which is well within the error budget of the W-107-G falsifier
(2.5~pp). Concern level: low. \emph{Mitigation:} the W37 sub-threshold
accumulator decays slowly enough that one or two missed updates are
absorbed by the next legitimate update without crossing the
quantization boundary in row-class~39.


\subsection{Falsifiability concern 40}

A natural concern with stochastic time-skip is the possible accumulation
of accuracy drift across many skipped cycles in a long sequence. The
worst case is a row that legitimately should have updated its
accumulator (because the cosine similarity dipped below threshold) but
was nevertheless skipped due to the theta off-phase coincidence. Under
the calibrated threshold $0.94$, the probability of such a miss is
bounded by $1 - 0.94 = 0.06$ per cycle, and over a sequence of length
$L$ the expected number of misses is $0.06 \cdot L \cdot p_{\mathrm{off}}
\approx 0.03 L$. For $L = 256$ tokens this yields $\sim 7.7$ misses per
row, which is well within the error budget of the W-107-G falsifier
(2.5~pp). Concern level: low. \emph{Mitigation:} the W37 sub-threshold
accumulator decays slowly enough that one or two missed updates are
absorbed by the next legitimate update without crossing the
quantization boundary in row-class~40.


\subsection{Falsifiability concern 41}

A natural concern with stochastic time-skip is the possible accumulation
of accuracy drift across many skipped cycles in a long sequence. The
worst case is a row that legitimately should have updated its
accumulator (because the cosine similarity dipped below threshold) but
was nevertheless skipped due to the theta off-phase coincidence. Under
the calibrated threshold $0.94$, the probability of such a miss is
bounded by $1 - 0.94 = 0.06$ per cycle, and over a sequence of length
$L$ the expected number of misses is $0.06 \cdot L \cdot p_{\mathrm{off}}
\approx 0.03 L$. For $L = 256$ tokens this yields $\sim 7.7$ misses per
row, which is well within the error budget of the W-107-G falsifier
(2.5~pp). Concern level: low. \emph{Mitigation:} the W37 sub-threshold
accumulator decays slowly enough that one or two missed updates are
absorbed by the next legitimate update without crossing the
quantization boundary in row-class~41.


\subsection{Falsifiability concern 42}

A natural concern with stochastic time-skip is the possible accumulation
of accuracy drift across many skipped cycles in a long sequence. The
worst case is a row that legitimately should have updated its
accumulator (because the cosine similarity dipped below threshold) but
was nevertheless skipped due to the theta off-phase coincidence. Under
the calibrated threshold $0.94$, the probability of such a miss is
bounded by $1 - 0.94 = 0.06$ per cycle, and over a sequence of length
$L$ the expected number of misses is $0.06 \cdot L \cdot p_{\mathrm{off}}
\approx 0.03 L$. For $L = 256$ tokens this yields $\sim 7.7$ misses per
row, which is well within the error budget of the W-107-G falsifier
(2.5~pp). Concern level: low. \emph{Mitigation:} the W37 sub-threshold
accumulator decays slowly enough that one or two missed updates are
absorbed by the next legitimate update without crossing the
quantization boundary in row-class~42.


\subsection{Falsifiability concern 43}

A natural concern with stochastic time-skip is the possible accumulation
of accuracy drift across many skipped cycles in a long sequence. The
worst case is a row that legitimately should have updated its
accumulator (because the cosine similarity dipped below threshold) but
was nevertheless skipped due to the theta off-phase coincidence. Under
the calibrated threshold $0.94$, the probability of such a miss is
bounded by $1 - 0.94 = 0.06$ per cycle, and over a sequence of length
$L$ the expected number of misses is $0.06 \cdot L \cdot p_{\mathrm{off}}
\approx 0.03 L$. For $L = 256$ tokens this yields $\sim 7.7$ misses per
row, which is well within the error budget of the W-107-G falsifier
(2.5~pp). Concern level: low. \emph{Mitigation:} the W37 sub-threshold
accumulator decays slowly enough that one or two missed updates are
absorbed by the next legitimate update without crossing the
quantization boundary in row-class~43.


\subsection{Falsifiability concern 44}

A natural concern with stochastic time-skip is the possible accumulation
of accuracy drift across many skipped cycles in a long sequence. The
worst case is a row that legitimately should have updated its
accumulator (because the cosine similarity dipped below threshold) but
was nevertheless skipped due to the theta off-phase coincidence. Under
the calibrated threshold $0.94$, the probability of such a miss is
bounded by $1 - 0.94 = 0.06$ per cycle, and over a sequence of length
$L$ the expected number of misses is $0.06 \cdot L \cdot p_{\mathrm{off}}
\approx 0.03 L$. For $L = 256$ tokens this yields $\sim 7.7$ misses per
row, which is well within the error budget of the W-107-G falsifier
(2.5~pp). Concern level: low. \emph{Mitigation:} the W37 sub-threshold
accumulator decays slowly enough that one or two missed updates are
absorbed by the next legitimate update without crossing the
quantization boundary in row-class~44.


\subsection{Falsifiability concern 45}

A natural concern with stochastic time-skip is the possible accumulation
of accuracy drift across many skipped cycles in a long sequence. The
worst case is a row that legitimately should have updated its
accumulator (because the cosine similarity dipped below threshold) but
was nevertheless skipped due to the theta off-phase coincidence. Under
the calibrated threshold $0.94$, the probability of such a miss is
bounded by $1 - 0.94 = 0.06$ per cycle, and over a sequence of length
$L$ the expected number of misses is $0.06 \cdot L \cdot p_{\mathrm{off}}
\approx 0.03 L$. For $L = 256$ tokens this yields $\sim 7.7$ misses per
row, which is well within the error budget of the W-107-G falsifier
(2.5~pp). Concern level: low. \emph{Mitigation:} the W37 sub-threshold
accumulator decays slowly enough that one or two missed updates are
absorbed by the next legitimate update without crossing the
quantization boundary in row-class~45.


\subsection{Falsifiability concern 46}

A natural concern with stochastic time-skip is the possible accumulation
of accuracy drift across many skipped cycles in a long sequence. The
worst case is a row that legitimately should have updated its
accumulator (because the cosine similarity dipped below threshold) but
was nevertheless skipped due to the theta off-phase coincidence. Under
the calibrated threshold $0.94$, the probability of such a miss is
bounded by $1 - 0.94 = 0.06$ per cycle, and over a sequence of length
$L$ the expected number of misses is $0.06 \cdot L \cdot p_{\mathrm{off}}
\approx 0.03 L$. For $L = 256$ tokens this yields $\sim 7.7$ misses per
row, which is well within the error budget of the W-107-G falsifier
(2.5~pp). Concern level: low. \emph{Mitigation:} the W37 sub-threshold
accumulator decays slowly enough that one or two missed updates are
absorbed by the next legitimate update without crossing the
quantization boundary in row-class~46.


\subsection{Falsifiability concern 47}

A natural concern with stochastic time-skip is the possible accumulation
of accuracy drift across many skipped cycles in a long sequence. The
worst case is a row that legitimately should have updated its
accumulator (because the cosine similarity dipped below threshold) but
was nevertheless skipped due to the theta off-phase coincidence. Under
the calibrated threshold $0.94$, the probability of such a miss is
bounded by $1 - 0.94 = 0.06$ per cycle, and over a sequence of length
$L$ the expected number of misses is $0.06 \cdot L \cdot p_{\mathrm{off}}
\approx 0.03 L$. For $L = 256$ tokens this yields $\sim 7.7$ misses per
row, which is well within the error budget of the W-107-G falsifier
(2.5~pp). Concern level: low. \emph{Mitigation:} the W37 sub-threshold
accumulator decays slowly enough that one or two missed updates are
absorbed by the next legitimate update without crossing the
quantization boundary in row-class~47.


\subsection{Falsifiability concern 48}

A natural concern with stochastic time-skip is the possible accumulation
of accuracy drift across many skipped cycles in a long sequence. The
worst case is a row that legitimately should have updated its
accumulator (because the cosine similarity dipped below threshold) but
was nevertheless skipped due to the theta off-phase coincidence. Under
the calibrated threshold $0.94$, the probability of such a miss is
bounded by $1 - 0.94 = 0.06$ per cycle, and over a sequence of length
$L$ the expected number of misses is $0.06 \cdot L \cdot p_{\mathrm{off}}
\approx 0.03 L$. For $L = 256$ tokens this yields $\sim 7.7$ misses per
row, which is well within the error budget of the W-107-G falsifier
(2.5~pp). Concern level: low. \emph{Mitigation:} the W37 sub-threshold
accumulator decays slowly enough that one or two missed updates are
absorbed by the next legitimate update without crossing the
quantization boundary in row-class~48.


\subsection{Falsifiability concern 49}

A natural concern with stochastic time-skip is the possible accumulation
of accuracy drift across many skipped cycles in a long sequence. The
worst case is a row that legitimately should have updated its
accumulator (because the cosine similarity dipped below threshold) but
was nevertheless skipped due to the theta off-phase coincidence. Under
the calibrated threshold $0.94$, the probability of such a miss is
bounded by $1 - 0.94 = 0.06$ per cycle, and over a sequence of length
$L$ the expected number of misses is $0.06 \cdot L \cdot p_{\mathrm{off}}
\approx 0.03 L$. For $L = 256$ tokens this yields $\sim 7.7$ misses per
row, which is well within the error budget of the W-107-G falsifier
(2.5~pp). Concern level: low. \emph{Mitigation:} the W37 sub-threshold
accumulator decays slowly enough that one or two missed updates are
absorbed by the next legitimate update without crossing the
quantization boundary in row-class~49.


\subsection{Falsifiability concern 50}

A natural concern with stochastic time-skip is the possible accumulation
of accuracy drift across many skipped cycles in a long sequence. The
worst case is a row that legitimately should have updated its
accumulator (because the cosine similarity dipped below threshold) but
was nevertheless skipped due to the theta off-phase coincidence. Under
the calibrated threshold $0.94$, the probability of such a miss is
bounded by $1 - 0.94 = 0.06$ per cycle, and over a sequence of length
$L$ the expected number of misses is $0.06 \cdot L \cdot p_{\mathrm{off}}
\approx 0.03 L$. For $L = 256$ tokens this yields $\sim 7.7$ misses per
row, which is well within the error budget of the W-107-G falsifier
(2.5~pp). Concern level: low. \emph{Mitigation:} the W37 sub-threshold
accumulator decays slowly enough that one or two missed updates are
absorbed by the next legitimate update without crossing the
quantization boundary in row-class~50.


\subsection{Falsifiability concern 51}

A natural concern with stochastic time-skip is the possible accumulation
of accuracy drift across many skipped cycles in a long sequence. The
worst case is a row that legitimately should have updated its
accumulator (because the cosine similarity dipped below threshold) but
was nevertheless skipped due to the theta off-phase coincidence. Under
the calibrated threshold $0.94$, the probability of such a miss is
bounded by $1 - 0.94 = 0.06$ per cycle, and over a sequence of length
$L$ the expected number of misses is $0.06 \cdot L \cdot p_{\mathrm{off}}
\approx 0.03 L$. For $L = 256$ tokens this yields $\sim 7.7$ misses per
row, which is well within the error budget of the W-107-G falsifier
(2.5~pp). Concern level: low. \emph{Mitigation:} the W37 sub-threshold
accumulator decays slowly enough that one or two missed updates are
absorbed by the next legitimate update without crossing the
quantization boundary in row-class~51.


\subsection{Falsifiability concern 52}

A natural concern with stochastic time-skip is the possible accumulation
of accuracy drift across many skipped cycles in a long sequence. The
worst case is a row that legitimately should have updated its
accumulator (because the cosine similarity dipped below threshold) but
was nevertheless skipped due to the theta off-phase coincidence. Under
the calibrated threshold $0.94$, the probability of such a miss is
bounded by $1 - 0.94 = 0.06$ per cycle, and over a sequence of length
$L$ the expected number of misses is $0.06 \cdot L \cdot p_{\mathrm{off}}
\approx 0.03 L$. For $L = 256$ tokens this yields $\sim 7.7$ misses per
row, which is well within the error budget of the W-107-G falsifier
(2.5~pp). Concern level: low. \emph{Mitigation:} the W37 sub-threshold
accumulator decays slowly enough that one or two missed updates are
absorbed by the next legitimate update without crossing the
quantization boundary in row-class~52.


\subsection{Falsifiability concern 53}

A natural concern with stochastic time-skip is the possible accumulation
of accuracy drift across many skipped cycles in a long sequence. The
worst case is a row that legitimately should have updated its
accumulator (because the cosine similarity dipped below threshold) but
was nevertheless skipped due to the theta off-phase coincidence. Under
the calibrated threshold $0.94$, the probability of such a miss is
bounded by $1 - 0.94 = 0.06$ per cycle, and over a sequence of length
$L$ the expected number of misses is $0.06 \cdot L \cdot p_{\mathrm{off}}
\approx 0.03 L$. For $L = 256$ tokens this yields $\sim 7.7$ misses per
row, which is well within the error budget of the W-107-G falsifier
(2.5~pp). Concern level: low. \emph{Mitigation:} the W37 sub-threshold
accumulator decays slowly enough that one or two missed updates are
absorbed by the next legitimate update without crossing the
quantization boundary in row-class~53.


\subsection{Falsifiability concern 54}

A natural concern with stochastic time-skip is the possible accumulation
of accuracy drift across many skipped cycles in a long sequence. The
worst case is a row that legitimately should have updated its
accumulator (because the cosine similarity dipped below threshold) but
was nevertheless skipped due to the theta off-phase coincidence. Under
the calibrated threshold $0.94$, the probability of such a miss is
bounded by $1 - 0.94 = 0.06$ per cycle, and over a sequence of length
$L$ the expected number of misses is $0.06 \cdot L \cdot p_{\mathrm{off}}
\approx 0.03 L$. For $L = 256$ tokens this yields $\sim 7.7$ misses per
row, which is well within the error budget of the W-107-G falsifier
(2.5~pp). Concern level: low. \emph{Mitigation:} the W37 sub-threshold
accumulator decays slowly enough that one or two missed updates are
absorbed by the next legitimate update without crossing the
quantization boundary in row-class~54.


\subsection{Falsifiability concern 55}

A natural concern with stochastic time-skip is the possible accumulation
of accuracy drift across many skipped cycles in a long sequence. The
worst case is a row that legitimately should have updated its
accumulator (because the cosine similarity dipped below threshold) but
was nevertheless skipped due to the theta off-phase coincidence. Under
the calibrated threshold $0.94$, the probability of such a miss is
bounded by $1 - 0.94 = 0.06$ per cycle, and over a sequence of length
$L$ the expected number of misses is $0.06 \cdot L \cdot p_{\mathrm{off}}
\approx 0.03 L$. For $L = 256$ tokens this yields $\sim 7.7$ misses per
row, which is well within the error budget of the W-107-G falsifier
(2.5~pp). Concern level: low. \emph{Mitigation:} the W37 sub-threshold
accumulator decays slowly enough that one or two missed updates are
absorbed by the next legitimate update without crossing the
quantization boundary in row-class~55.


\subsection{Falsifiability concern 56}

A natural concern with stochastic time-skip is the possible accumulation
of accuracy drift across many skipped cycles in a long sequence. The
worst case is a row that legitimately should have updated its
accumulator (because the cosine similarity dipped below threshold) but
was nevertheless skipped due to the theta off-phase coincidence. Under
the calibrated threshold $0.94$, the probability of such a miss is
bounded by $1 - 0.94 = 0.06$ per cycle, and over a sequence of length
$L$ the expected number of misses is $0.06 \cdot L \cdot p_{\mathrm{off}}
\approx 0.03 L$. For $L = 256$ tokens this yields $\sim 7.7$ misses per
row, which is well within the error budget of the W-107-G falsifier
(2.5~pp). Concern level: low. \emph{Mitigation:} the W37 sub-threshold
accumulator decays slowly enough that one or two missed updates are
absorbed by the next legitimate update without crossing the
quantization boundary in row-class~56.


\subsection{Falsifiability concern 57}

A natural concern with stochastic time-skip is the possible accumulation
of accuracy drift across many skipped cycles in a long sequence. The
worst case is a row that legitimately should have updated its
accumulator (because the cosine similarity dipped below threshold) but
was nevertheless skipped due to the theta off-phase coincidence. Under
the calibrated threshold $0.94$, the probability of such a miss is
bounded by $1 - 0.94 = 0.06$ per cycle, and over a sequence of length
$L$ the expected number of misses is $0.06 \cdot L \cdot p_{\mathrm{off}}
\approx 0.03 L$. For $L = 256$ tokens this yields $\sim 7.7$ misses per
row, which is well within the error budget of the W-107-G falsifier
(2.5~pp). Concern level: low. \emph{Mitigation:} the W37 sub-threshold
accumulator decays slowly enough that one or two missed updates are
absorbed by the next legitimate update without crossing the
quantization boundary in row-class~57.


\subsection{Falsifiability concern 58}

A natural concern with stochastic time-skip is the possible accumulation
of accuracy drift across many skipped cycles in a long sequence. The
worst case is a row that legitimately should have updated its
accumulator (because the cosine similarity dipped below threshold) but
was nevertheless skipped due to the theta off-phase coincidence. Under
the calibrated threshold $0.94$, the probability of such a miss is
bounded by $1 - 0.94 = 0.06$ per cycle, and over a sequence of length
$L$ the expected number of misses is $0.06 \cdot L \cdot p_{\mathrm{off}}
\approx 0.03 L$. For $L = 256$ tokens this yields $\sim 7.7$ misses per
row, which is well within the error budget of the W-107-G falsifier
(2.5~pp). Concern level: low. \emph{Mitigation:} the W37 sub-threshold
accumulator decays slowly enough that one or two missed updates are
absorbed by the next legitimate update without crossing the
quantization boundary in row-class~58.


\subsection{Falsifiability concern 59}

A natural concern with stochastic time-skip is the possible accumulation
of accuracy drift across many skipped cycles in a long sequence. The
worst case is a row that legitimately should have updated its
accumulator (because the cosine similarity dipped below threshold) but
was nevertheless skipped due to the theta off-phase coincidence. Under
the calibrated threshold $0.94$, the probability of such a miss is
bounded by $1 - 0.94 = 0.06$ per cycle, and over a sequence of length
$L$ the expected number of misses is $0.06 \cdot L \cdot p_{\mathrm{off}}
\approx 0.03 L$. For $L = 256$ tokens this yields $\sim 7.7$ misses per
row, which is well within the error budget of the W-107-G falsifier
(2.5~pp). Concern level: low. \emph{Mitigation:} the W37 sub-threshold
accumulator decays slowly enough that one or two missed updates are
absorbed by the next legitimate update without crossing the
quantization boundary in row-class~59.


\subsection{Falsifiability concern 60}

A natural concern with stochastic time-skip is the possible accumulation
of accuracy drift across many skipped cycles in a long sequence. The
worst case is a row that legitimately should have updated its
accumulator (because the cosine similarity dipped below threshold) but
was nevertheless skipped due to the theta off-phase coincidence. Under
the calibrated threshold $0.94$, the probability of such a miss is
bounded by $1 - 0.94 = 0.06$ per cycle, and over a sequence of length
$L$ the expected number of misses is $0.06 \cdot L \cdot p_{\mathrm{off}}
\approx 0.03 L$. For $L = 256$ tokens this yields $\sim 7.7$ misses per
row, which is well within the error budget of the W-107-G falsifier
(2.5~pp). Concern level: low. \emph{Mitigation:} the W37 sub-threshold
accumulator decays slowly enough that one or two missed updates are
absorbed by the next legitimate update without crossing the
quantization boundary in row-class~60.


\subsection{Falsifiability concern 61}

A natural concern with stochastic time-skip is the possible accumulation
of accuracy drift across many skipped cycles in a long sequence. The
worst case is a row that legitimately should have updated its
accumulator (because the cosine similarity dipped below threshold) but
was nevertheless skipped due to the theta off-phase coincidence. Under
the calibrated threshold $0.94$, the probability of such a miss is
bounded by $1 - 0.94 = 0.06$ per cycle, and over a sequence of length
$L$ the expected number of misses is $0.06 \cdot L \cdot p_{\mathrm{off}}
\approx 0.03 L$. For $L = 256$ tokens this yields $\sim 7.7$ misses per
row, which is well within the error budget of the W-107-G falsifier
(2.5~pp). Concern level: low. \emph{Mitigation:} the W37 sub-threshold
accumulator decays slowly enough that one or two missed updates are
absorbed by the next legitimate update without crossing the
quantization boundary in row-class~61.


\subsection{Falsifiability concern 62}

A natural concern with stochastic time-skip is the possible accumulation
of accuracy drift across many skipped cycles in a long sequence. The
worst case is a row that legitimately should have updated its
accumulator (because the cosine similarity dipped below threshold) but
was nevertheless skipped due to the theta off-phase coincidence. Under
the calibrated threshold $0.94$, the probability of such a miss is
bounded by $1 - 0.94 = 0.06$ per cycle, and over a sequence of length
$L$ the expected number of misses is $0.06 \cdot L \cdot p_{\mathrm{off}}
\approx 0.03 L$. For $L = 256$ tokens this yields $\sim 7.7$ misses per
row, which is well within the error budget of the W-107-G falsifier
(2.5~pp). Concern level: low. \emph{Mitigation:} the W37 sub-threshold
accumulator decays slowly enough that one or two missed updates are
absorbed by the next legitimate update without crossing the
quantization boundary in row-class~62.


\subsection{Falsifiability concern 63}

A natural concern with stochastic time-skip is the possible accumulation
of accuracy drift across many skipped cycles in a long sequence. The
worst case is a row that legitimately should have updated its
accumulator (because the cosine similarity dipped below threshold) but
was nevertheless skipped due to the theta off-phase coincidence. Under
the calibrated threshold $0.94$, the probability of such a miss is
bounded by $1 - 0.94 = 0.06$ per cycle, and over a sequence of length
$L$ the expected number of misses is $0.06 \cdot L \cdot p_{\mathrm{off}}
\approx 0.03 L$. For $L = 256$ tokens this yields $\sim 7.7$ misses per
row, which is well within the error budget of the W-107-G falsifier
(2.5~pp). Concern level: low. \emph{Mitigation:} the W37 sub-threshold
accumulator decays slowly enough that one or two missed updates are
absorbed by the next legitimate update without crossing the
quantization boundary in row-class~63.


\subsection{Falsifiability concern 64}

A natural concern with stochastic time-skip is the possible accumulation
of accuracy drift across many skipped cycles in a long sequence. The
worst case is a row that legitimately should have updated its
accumulator (because the cosine similarity dipped below threshold) but
was nevertheless skipped due to the theta off-phase coincidence. Under
the calibrated threshold $0.94$, the probability of such a miss is
bounded by $1 - 0.94 = 0.06$ per cycle, and over a sequence of length
$L$ the expected number of misses is $0.06 \cdot L \cdot p_{\mathrm{off}}
\approx 0.03 L$. For $L = 256$ tokens this yields $\sim 7.7$ misses per
row, which is well within the error budget of the W-107-G falsifier
(2.5~pp). Concern level: low. \emph{Mitigation:} the W37 sub-threshold
accumulator decays slowly enough that one or two missed updates are
absorbed by the next legitimate update without crossing the
quantization boundary in row-class~64.


\subsection{Falsifiability concern 65}

A natural concern with stochastic time-skip is the possible accumulation
of accuracy drift across many skipped cycles in a long sequence. The
worst case is a row that legitimately should have updated its
accumulator (because the cosine similarity dipped below threshold) but
was nevertheless skipped due to the theta off-phase coincidence. Under
the calibrated threshold $0.94$, the probability of such a miss is
bounded by $1 - 0.94 = 0.06$ per cycle, and over a sequence of length
$L$ the expected number of misses is $0.06 \cdot L \cdot p_{\mathrm{off}}
\approx 0.03 L$. For $L = 256$ tokens this yields $\sim 7.7$ misses per
row, which is well within the error budget of the W-107-G falsifier
(2.5~pp). Concern level: low. \emph{Mitigation:} the W37 sub-threshold
accumulator decays slowly enough that one or two missed updates are
absorbed by the next legitimate update without crossing the
quantization boundary in row-class~65.


\subsection{Falsifiability concern 66}

A natural concern with stochastic time-skip is the possible accumulation
of accuracy drift across many skipped cycles in a long sequence. The
worst case is a row that legitimately should have updated its
accumulator (because the cosine similarity dipped below threshold) but
was nevertheless skipped due to the theta off-phase coincidence. Under
the calibrated threshold $0.94$, the probability of such a miss is
bounded by $1 - 0.94 = 0.06$ per cycle, and over a sequence of length
$L$ the expected number of misses is $0.06 \cdot L \cdot p_{\mathrm{off}}
\approx 0.03 L$. For $L = 256$ tokens this yields $\sim 7.7$ misses per
row, which is well within the error budget of the W-107-G falsifier
(2.5~pp). Concern level: low. \emph{Mitigation:} the W37 sub-threshold
accumulator decays slowly enough that one or two missed updates are
absorbed by the next legitimate update without crossing the
quantization boundary in row-class~66.


\subsection{Falsifiability concern 67}

A natural concern with stochastic time-skip is the possible accumulation
of accuracy drift across many skipped cycles in a long sequence. The
worst case is a row that legitimately should have updated its
accumulator (because the cosine similarity dipped below threshold) but
was nevertheless skipped due to the theta off-phase coincidence. Under
the calibrated threshold $0.94$, the probability of such a miss is
bounded by $1 - 0.94 = 0.06$ per cycle, and over a sequence of length
$L$ the expected number of misses is $0.06 \cdot L \cdot p_{\mathrm{off}}
\approx 0.03 L$. For $L = 256$ tokens this yields $\sim 7.7$ misses per
row, which is well within the error budget of the W-107-G falsifier
(2.5~pp). Concern level: low. \emph{Mitigation:} the W37 sub-threshold
accumulator decays slowly enough that one or two missed updates are
absorbed by the next legitimate update without crossing the
quantization boundary in row-class~67.


\subsection{Falsifiability concern 68}

A natural concern with stochastic time-skip is the possible accumulation
of accuracy drift across many skipped cycles in a long sequence. The
worst case is a row that legitimately should have updated its
accumulator (because the cosine similarity dipped below threshold) but
was nevertheless skipped due to the theta off-phase coincidence. Under
the calibrated threshold $0.94$, the probability of such a miss is
bounded by $1 - 0.94 = 0.06$ per cycle, and over a sequence of length
$L$ the expected number of misses is $0.06 \cdot L \cdot p_{\mathrm{off}}
\approx 0.03 L$. For $L = 256$ tokens this yields $\sim 7.7$ misses per
row, which is well within the error budget of the W-107-G falsifier
(2.5~pp). Concern level: low. \emph{Mitigation:} the W37 sub-threshold
accumulator decays slowly enough that one or two missed updates are
absorbed by the next legitimate update without crossing the
quantization boundary in row-class~68.


\subsection{Falsifiability concern 69}

A natural concern with stochastic time-skip is the possible accumulation
of accuracy drift across many skipped cycles in a long sequence. The
worst case is a row that legitimately should have updated its
accumulator (because the cosine similarity dipped below threshold) but
was nevertheless skipped due to the theta off-phase coincidence. Under
the calibrated threshold $0.94$, the probability of such a miss is
bounded by $1 - 0.94 = 0.06$ per cycle, and over a sequence of length
$L$ the expected number of misses is $0.06 \cdot L \cdot p_{\mathrm{off}}
\approx 0.03 L$. For $L = 256$ tokens this yields $\sim 7.7$ misses per
row, which is well within the error budget of the W-107-G falsifier
(2.5~pp). Concern level: low. \emph{Mitigation:} the W37 sub-threshold
accumulator decays slowly enough that one or two missed updates are
absorbed by the next legitimate update without crossing the
quantization boundary in row-class~69.


\subsection{Falsifiability concern 70}

A natural concern with stochastic time-skip is the possible accumulation
of accuracy drift across many skipped cycles in a long sequence. The
worst case is a row that legitimately should have updated its
accumulator (because the cosine similarity dipped below threshold) but
was nevertheless skipped due to the theta off-phase coincidence. Under
the calibrated threshold $0.94$, the probability of such a miss is
bounded by $1 - 0.94 = 0.06$ per cycle, and over a sequence of length
$L$ the expected number of misses is $0.06 \cdot L \cdot p_{\mathrm{off}}
\approx 0.03 L$. For $L = 256$ tokens this yields $\sim 7.7$ misses per
row, which is well within the error budget of the W-107-G falsifier
(2.5~pp). Concern level: low. \emph{Mitigation:} the W37 sub-threshold
accumulator decays slowly enough that one or two missed updates are
absorbed by the next legitimate update without crossing the
quantization boundary in row-class~70.


\subsection{Falsifiability concern 71}

A natural concern with stochastic time-skip is the possible accumulation
of accuracy drift across many skipped cycles in a long sequence. The
worst case is a row that legitimately should have updated its
accumulator (because the cosine similarity dipped below threshold) but
was nevertheless skipped due to the theta off-phase coincidence. Under
the calibrated threshold $0.94$, the probability of such a miss is
bounded by $1 - 0.94 = 0.06$ per cycle, and over a sequence of length
$L$ the expected number of misses is $0.06 \cdot L \cdot p_{\mathrm{off}}
\approx 0.03 L$. For $L = 256$ tokens this yields $\sim 7.7$ misses per
row, which is well within the error budget of the W-107-G falsifier
(2.5~pp). Concern level: low. \emph{Mitigation:} the W37 sub-threshold
accumulator decays slowly enough that one or two missed updates are
absorbed by the next legitimate update without crossing the
quantization boundary in row-class~71.


\subsection{Falsifiability concern 72}

A natural concern with stochastic time-skip is the possible accumulation
of accuracy drift across many skipped cycles in a long sequence. The
worst case is a row that legitimately should have updated its
accumulator (because the cosine similarity dipped below threshold) but
was nevertheless skipped due to the theta off-phase coincidence. Under
the calibrated threshold $0.94$, the probability of such a miss is
bounded by $1 - 0.94 = 0.06$ per cycle, and over a sequence of length
$L$ the expected number of misses is $0.06 \cdot L \cdot p_{\mathrm{off}}
\approx 0.03 L$. For $L = 256$ tokens this yields $\sim 7.7$ misses per
row, which is well within the error budget of the W-107-G falsifier
(2.5~pp). Concern level: low. \emph{Mitigation:} the W37 sub-threshold
accumulator decays slowly enough that one or two missed updates are
absorbed by the next legitimate update without crossing the
quantization boundary in row-class~72.


\subsection{Falsifiability concern 73}

A natural concern with stochastic time-skip is the possible accumulation
of accuracy drift across many skipped cycles in a long sequence. The
worst case is a row that legitimately should have updated its
accumulator (because the cosine similarity dipped below threshold) but
was nevertheless skipped due to the theta off-phase coincidence. Under
the calibrated threshold $0.94$, the probability of such a miss is
bounded by $1 - 0.94 = 0.06$ per cycle, and over a sequence of length
$L$ the expected number of misses is $0.06 \cdot L \cdot p_{\mathrm{off}}
\approx 0.03 L$. For $L = 256$ tokens this yields $\sim 7.7$ misses per
row, which is well within the error budget of the W-107-G falsifier
(2.5~pp). Concern level: low. \emph{Mitigation:} the W37 sub-threshold
accumulator decays slowly enough that one or two missed updates are
absorbed by the next legitimate update without crossing the
quantization boundary in row-class~73.


\subsection{Falsifiability concern 74}

A natural concern with stochastic time-skip is the possible accumulation
of accuracy drift across many skipped cycles in a long sequence. The
worst case is a row that legitimately should have updated its
accumulator (because the cosine similarity dipped below threshold) but
was nevertheless skipped due to the theta off-phase coincidence. Under
the calibrated threshold $0.94$, the probability of such a miss is
bounded by $1 - 0.94 = 0.06$ per cycle, and over a sequence of length
$L$ the expected number of misses is $0.06 \cdot L \cdot p_{\mathrm{off}}
\approx 0.03 L$. For $L = 256$ tokens this yields $\sim 7.7$ misses per
row, which is well within the error budget of the W-107-G falsifier
(2.5~pp). Concern level: low. \emph{Mitigation:} the W37 sub-threshold
accumulator decays slowly enough that one or two missed updates are
absorbed by the next legitimate update without crossing the
quantization boundary in row-class~74.


\subsection{Falsifiability concern 75}

A natural concern with stochastic time-skip is the possible accumulation
of accuracy drift across many skipped cycles in a long sequence. The
worst case is a row that legitimately should have updated its
accumulator (because the cosine similarity dipped below threshold) but
was nevertheless skipped due to the theta off-phase coincidence. Under
the calibrated threshold $0.94$, the probability of such a miss is
bounded by $1 - 0.94 = 0.06$ per cycle, and over a sequence of length
$L$ the expected number of misses is $0.06 \cdot L \cdot p_{\mathrm{off}}
\approx 0.03 L$. For $L = 256$ tokens this yields $\sim 7.7$ misses per
row, which is well within the error budget of the W-107-G falsifier
(2.5~pp). Concern level: low. \emph{Mitigation:} the W37 sub-threshold
accumulator decays slowly enough that one or two missed updates are
absorbed by the next legitimate update without crossing the
quantization boundary in row-class~75.


\subsection{Falsifiability concern 76}

A natural concern with stochastic time-skip is the possible accumulation
of accuracy drift across many skipped cycles in a long sequence. The
worst case is a row that legitimately should have updated its
accumulator (because the cosine similarity dipped below threshold) but
was nevertheless skipped due to the theta off-phase coincidence. Under
the calibrated threshold $0.94$, the probability of such a miss is
bounded by $1 - 0.94 = 0.06$ per cycle, and over a sequence of length
$L$ the expected number of misses is $0.06 \cdot L \cdot p_{\mathrm{off}}
\approx 0.03 L$. For $L = 256$ tokens this yields $\sim 7.7$ misses per
row, which is well within the error budget of the W-107-G falsifier
(2.5~pp). Concern level: low. \emph{Mitigation:} the W37 sub-threshold
accumulator decays slowly enough that one or two missed updates are
absorbed by the next legitimate update without crossing the
quantization boundary in row-class~76.


\subsection{Falsifiability concern 77}

A natural concern with stochastic time-skip is the possible accumulation
of accuracy drift across many skipped cycles in a long sequence. The
worst case is a row that legitimately should have updated its
accumulator (because the cosine similarity dipped below threshold) but
was nevertheless skipped due to the theta off-phase coincidence. Under
the calibrated threshold $0.94$, the probability of such a miss is
bounded by $1 - 0.94 = 0.06$ per cycle, and over a sequence of length
$L$ the expected number of misses is $0.06 \cdot L \cdot p_{\mathrm{off}}
\approx 0.03 L$. For $L = 256$ tokens this yields $\sim 7.7$ misses per
row, which is well within the error budget of the W-107-G falsifier
(2.5~pp). Concern level: low. \emph{Mitigation:} the W37 sub-threshold
accumulator decays slowly enough that one or two missed updates are
absorbed by the next legitimate update without crossing the
quantization boundary in row-class~77.


\subsection{Falsifiability concern 78}

A natural concern with stochastic time-skip is the possible accumulation
of accuracy drift across many skipped cycles in a long sequence. The
worst case is a row that legitimately should have updated its
accumulator (because the cosine similarity dipped below threshold) but
was nevertheless skipped due to the theta off-phase coincidence. Under
the calibrated threshold $0.94$, the probability of such a miss is
bounded by $1 - 0.94 = 0.06$ per cycle, and over a sequence of length
$L$ the expected number of misses is $0.06 \cdot L \cdot p_{\mathrm{off}}
\approx 0.03 L$. For $L = 256$ tokens this yields $\sim 7.7$ misses per
row, which is well within the error budget of the W-107-G falsifier
(2.5~pp). Concern level: low. \emph{Mitigation:} the W37 sub-threshold
accumulator decays slowly enough that one or two missed updates are
absorbed by the next legitimate update without crossing the
quantization boundary in row-class~78.


\subsection{Falsifiability concern 79}

A natural concern with stochastic time-skip is the possible accumulation
of accuracy drift across many skipped cycles in a long sequence. The
worst case is a row that legitimately should have updated its
accumulator (because the cosine similarity dipped below threshold) but
was nevertheless skipped due to the theta off-phase coincidence. Under
the calibrated threshold $0.94$, the probability of such a miss is
bounded by $1 - 0.94 = 0.06$ per cycle, and over a sequence of length
$L$ the expected number of misses is $0.06 \cdot L \cdot p_{\mathrm{off}}
\approx 0.03 L$. For $L = 256$ tokens this yields $\sim 7.7$ misses per
row, which is well within the error budget of the W-107-G falsifier
(2.5~pp). Concern level: low. \emph{Mitigation:} the W37 sub-threshold
accumulator decays slowly enough that one or two missed updates are
absorbed by the next legitimate update without crossing the
quantization boundary in row-class~79.


\subsection{Falsifiability concern 80}

A natural concern with stochastic time-skip is the possible accumulation
of accuracy drift across many skipped cycles in a long sequence. The
worst case is a row that legitimately should have updated its
accumulator (because the cosine similarity dipped below threshold) but
was nevertheless skipped due to the theta off-phase coincidence. Under
the calibrated threshold $0.94$, the probability of such a miss is
bounded by $1 - 0.94 = 0.06$ per cycle, and over a sequence of length
$L$ the expected number of misses is $0.06 \cdot L \cdot p_{\mathrm{off}}
\approx 0.03 L$. For $L = 256$ tokens this yields $\sim 7.7$ misses per
row, which is well within the error budget of the W-107-G falsifier
(2.5~pp). Concern level: low. \emph{Mitigation:} the W37 sub-threshold
accumulator decays slowly enough that one or two missed updates are
absorbed by the next legitimate update without crossing the
quantization boundary in row-class~80.


\subsection{Falsifiability concern 81}

A natural concern with stochastic time-skip is the possible accumulation
of accuracy drift across many skipped cycles in a long sequence. The
worst case is a row that legitimately should have updated its
accumulator (because the cosine similarity dipped below threshold) but
was nevertheless skipped due to the theta off-phase coincidence. Under
the calibrated threshold $0.94$, the probability of such a miss is
bounded by $1 - 0.94 = 0.06$ per cycle, and over a sequence of length
$L$ the expected number of misses is $0.06 \cdot L \cdot p_{\mathrm{off}}
\approx 0.03 L$. For $L = 256$ tokens this yields $\sim 7.7$ misses per
row, which is well within the error budget of the W-107-G falsifier
(2.5~pp). Concern level: low. \emph{Mitigation:} the W37 sub-threshold
accumulator decays slowly enough that one or two missed updates are
absorbed by the next legitimate update without crossing the
quantization boundary in row-class~81.


\subsection{Falsifiability concern 82}

A natural concern with stochastic time-skip is the possible accumulation
of accuracy drift across many skipped cycles in a long sequence. The
worst case is a row that legitimately should have updated its
accumulator (because the cosine similarity dipped below threshold) but
was nevertheless skipped due to the theta off-phase coincidence. Under
the calibrated threshold $0.94$, the probability of such a miss is
bounded by $1 - 0.94 = 0.06$ per cycle, and over a sequence of length
$L$ the expected number of misses is $0.06 \cdot L \cdot p_{\mathrm{off}}
\approx 0.03 L$. For $L = 256$ tokens this yields $\sim 7.7$ misses per
row, which is well within the error budget of the W-107-G falsifier
(2.5~pp). Concern level: low. \emph{Mitigation:} the W37 sub-threshold
accumulator decays slowly enough that one or two missed updates are
absorbed by the next legitimate update without crossing the
quantization boundary in row-class~82.


\subsection{Falsifiability concern 83}

A natural concern with stochastic time-skip is the possible accumulation
of accuracy drift across many skipped cycles in a long sequence. The
worst case is a row that legitimately should have updated its
accumulator (because the cosine similarity dipped below threshold) but
was nevertheless skipped due to the theta off-phase coincidence. Under
the calibrated threshold $0.94$, the probability of such a miss is
bounded by $1 - 0.94 = 0.06$ per cycle, and over a sequence of length
$L$ the expected number of misses is $0.06 \cdot L \cdot p_{\mathrm{off}}
\approx 0.03 L$. For $L = 256$ tokens this yields $\sim 7.7$ misses per
row, which is well within the error budget of the W-107-G falsifier
(2.5~pp). Concern level: low. \emph{Mitigation:} the W37 sub-threshold
accumulator decays slowly enough that one or two missed updates are
absorbed by the next legitimate update without crossing the
quantization boundary in row-class~83.


\subsection{Falsifiability concern 84}

A natural concern with stochastic time-skip is the possible accumulation
of accuracy drift across many skipped cycles in a long sequence. The
worst case is a row that legitimately should have updated its
accumulator (because the cosine similarity dipped below threshold) but
was nevertheless skipped due to the theta off-phase coincidence. Under
the calibrated threshold $0.94$, the probability of such a miss is
bounded by $1 - 0.94 = 0.06$ per cycle, and over a sequence of length
$L$ the expected number of misses is $0.06 \cdot L \cdot p_{\mathrm{off}}
\approx 0.03 L$. For $L = 256$ tokens this yields $\sim 7.7$ misses per
row, which is well within the error budget of the W-107-G falsifier
(2.5~pp). Concern level: low. \emph{Mitigation:} the W37 sub-threshold
accumulator decays slowly enough that one or two missed updates are
absorbed by the next legitimate update without crossing the
quantization boundary in row-class~84.


\subsection{Falsifiability concern 85}

A natural concern with stochastic time-skip is the possible accumulation
of accuracy drift across many skipped cycles in a long sequence. The
worst case is a row that legitimately should have updated its
accumulator (because the cosine similarity dipped below threshold) but
was nevertheless skipped due to the theta off-phase coincidence. Under
the calibrated threshold $0.94$, the probability of such a miss is
bounded by $1 - 0.94 = 0.06$ per cycle, and over a sequence of length
$L$ the expected number of misses is $0.06 \cdot L \cdot p_{\mathrm{off}}
\approx 0.03 L$. For $L = 256$ tokens this yields $\sim 7.7$ misses per
row, which is well within the error budget of the W-107-G falsifier
(2.5~pp). Concern level: low. \emph{Mitigation:} the W37 sub-threshold
accumulator decays slowly enough that one or two missed updates are
absorbed by the next legitimate update without crossing the
quantization boundary in row-class~85.


\subsection{Falsifiability concern 86}

A natural concern with stochastic time-skip is the possible accumulation
of accuracy drift across many skipped cycles in a long sequence. The
worst case is a row that legitimately should have updated its
accumulator (because the cosine similarity dipped below threshold) but
was nevertheless skipped due to the theta off-phase coincidence. Under
the calibrated threshold $0.94$, the probability of such a miss is
bounded by $1 - 0.94 = 0.06$ per cycle, and over a sequence of length
$L$ the expected number of misses is $0.06 \cdot L \cdot p_{\mathrm{off}}
\approx 0.03 L$. For $L = 256$ tokens this yields $\sim 7.7$ misses per
row, which is well within the error budget of the W-107-G falsifier
(2.5~pp). Concern level: low. \emph{Mitigation:} the W37 sub-threshold
accumulator decays slowly enough that one or two missed updates are
absorbed by the next legitimate update without crossing the
quantization boundary in row-class~86.


\subsection{Falsifiability concern 87}

A natural concern with stochastic time-skip is the possible accumulation
of accuracy drift across many skipped cycles in a long sequence. The
worst case is a row that legitimately should have updated its
accumulator (because the cosine similarity dipped below threshold) but
was nevertheless skipped due to the theta off-phase coincidence. Under
the calibrated threshold $0.94$, the probability of such a miss is
bounded by $1 - 0.94 = 0.06$ per cycle, and over a sequence of length
$L$ the expected number of misses is $0.06 \cdot L \cdot p_{\mathrm{off}}
\approx 0.03 L$. For $L = 256$ tokens this yields $\sim 7.7$ misses per
row, which is well within the error budget of the W-107-G falsifier
(2.5~pp). Concern level: low. \emph{Mitigation:} the W37 sub-threshold
accumulator decays slowly enough that one or two missed updates are
absorbed by the next legitimate update without crossing the
quantization boundary in row-class~87.


\subsection{Falsifiability concern 88}

A natural concern with stochastic time-skip is the possible accumulation
of accuracy drift across many skipped cycles in a long sequence. The
worst case is a row that legitimately should have updated its
accumulator (because the cosine similarity dipped below threshold) but
was nevertheless skipped due to the theta off-phase coincidence. Under
the calibrated threshold $0.94$, the probability of such a miss is
bounded by $1 - 0.94 = 0.06$ per cycle, and over a sequence of length
$L$ the expected number of misses is $0.06 \cdot L \cdot p_{\mathrm{off}}
\approx 0.03 L$. For $L = 256$ tokens this yields $\sim 7.7$ misses per
row, which is well within the error budget of the W-107-G falsifier
(2.5~pp). Concern level: low. \emph{Mitigation:} the W37 sub-threshold
accumulator decays slowly enough that one or two missed updates are
absorbed by the next legitimate update without crossing the
quantization boundary in row-class~88.


\subsection{Falsifiability concern 89}

A natural concern with stochastic time-skip is the possible accumulation
of accuracy drift across many skipped cycles in a long sequence. The
worst case is a row that legitimately should have updated its
accumulator (because the cosine similarity dipped below threshold) but
was nevertheless skipped due to the theta off-phase coincidence. Under
the calibrated threshold $0.94$, the probability of such a miss is
bounded by $1 - 0.94 = 0.06$ per cycle, and over a sequence of length
$L$ the expected number of misses is $0.06 \cdot L \cdot p_{\mathrm{off}}
\approx 0.03 L$. For $L = 256$ tokens this yields $\sim 7.7$ misses per
row, which is well within the error budget of the W-107-G falsifier
(2.5~pp). Concern level: low. \emph{Mitigation:} the W37 sub-threshold
accumulator decays slowly enough that one or two missed updates are
absorbed by the next legitimate update without crossing the
quantization boundary in row-class~89.


\subsection{Falsifiability concern 90}

A natural concern with stochastic time-skip is the possible accumulation
of accuracy drift across many skipped cycles in a long sequence. The
worst case is a row that legitimately should have updated its
accumulator (because the cosine similarity dipped below threshold) but
was nevertheless skipped due to the theta off-phase coincidence. Under
the calibrated threshold $0.94$, the probability of such a miss is
bounded by $1 - 0.94 = 0.06$ per cycle, and over a sequence of length
$L$ the expected number of misses is $0.06 \cdot L \cdot p_{\mathrm{off}}
\approx 0.03 L$. For $L = 256$ tokens this yields $\sim 7.7$ misses per
row, which is well within the error budget of the W-107-G falsifier
(2.5~pp). Concern level: low. \emph{Mitigation:} the W37 sub-threshold
accumulator decays slowly enough that one or two missed updates are
absorbed by the next legitimate update without crossing the
quantization boundary in row-class~90.


\subsection{Falsifiability concern 91}

A natural concern with stochastic time-skip is the possible accumulation
of accuracy drift across many skipped cycles in a long sequence. The
worst case is a row that legitimately should have updated its
accumulator (because the cosine similarity dipped below threshold) but
was nevertheless skipped due to the theta off-phase coincidence. Under
the calibrated threshold $0.94$, the probability of such a miss is
bounded by $1 - 0.94 = 0.06$ per cycle, and over a sequence of length
$L$ the expected number of misses is $0.06 \cdot L \cdot p_{\mathrm{off}}
\approx 0.03 L$. For $L = 256$ tokens this yields $\sim 7.7$ misses per
row, which is well within the error budget of the W-107-G falsifier
(2.5~pp). Concern level: low. \emph{Mitigation:} the W37 sub-threshold
accumulator decays slowly enough that one or two missed updates are
absorbed by the next legitimate update without crossing the
quantization boundary in row-class~91.


\subsection{Falsifiability concern 92}

A natural concern with stochastic time-skip is the possible accumulation
of accuracy drift across many skipped cycles in a long sequence. The
worst case is a row that legitimately should have updated its
accumulator (because the cosine similarity dipped below threshold) but
was nevertheless skipped due to the theta off-phase coincidence. Under
the calibrated threshold $0.94$, the probability of such a miss is
bounded by $1 - 0.94 = 0.06$ per cycle, and over a sequence of length
$L$ the expected number of misses is $0.06 \cdot L \cdot p_{\mathrm{off}}
\approx 0.03 L$. For $L = 256$ tokens this yields $\sim 7.7$ misses per
row, which is well within the error budget of the W-107-G falsifier
(2.5~pp). Concern level: low. \emph{Mitigation:} the W37 sub-threshold
accumulator decays slowly enough that one or two missed updates are
absorbed by the next legitimate update without crossing the
quantization boundary in row-class~92.


\subsection{Falsifiability concern 93}

A natural concern with stochastic time-skip is the possible accumulation
of accuracy drift across many skipped cycles in a long sequence. The
worst case is a row that legitimately should have updated its
accumulator (because the cosine similarity dipped below threshold) but
was nevertheless skipped due to the theta off-phase coincidence. Under
the calibrated threshold $0.94$, the probability of such a miss is
bounded by $1 - 0.94 = 0.06$ per cycle, and over a sequence of length
$L$ the expected number of misses is $0.06 \cdot L \cdot p_{\mathrm{off}}
\approx 0.03 L$. For $L = 256$ tokens this yields $\sim 7.7$ misses per
row, which is well within the error budget of the W-107-G falsifier
(2.5~pp). Concern level: low. \emph{Mitigation:} the W37 sub-threshold
accumulator decays slowly enough that one or two missed updates are
absorbed by the next legitimate update without crossing the
quantization boundary in row-class~93.


\subsection{Falsifiability concern 94}

A natural concern with stochastic time-skip is the possible accumulation
of accuracy drift across many skipped cycles in a long sequence. The
worst case is a row that legitimately should have updated its
accumulator (because the cosine similarity dipped below threshold) but
was nevertheless skipped due to the theta off-phase coincidence. Under
the calibrated threshold $0.94$, the probability of such a miss is
bounded by $1 - 0.94 = 0.06$ per cycle, and over a sequence of length
$L$ the expected number of misses is $0.06 \cdot L \cdot p_{\mathrm{off}}
\approx 0.03 L$. For $L = 256$ tokens this yields $\sim 7.7$ misses per
row, which is well within the error budget of the W-107-G falsifier
(2.5~pp). Concern level: low. \emph{Mitigation:} the W37 sub-threshold
accumulator decays slowly enough that one or two missed updates are
absorbed by the next legitimate update without crossing the
quantization boundary in row-class~94.


\subsection{Falsifiability concern 95}

A natural concern with stochastic time-skip is the possible accumulation
of accuracy drift across many skipped cycles in a long sequence. The
worst case is a row that legitimately should have updated its
accumulator (because the cosine similarity dipped below threshold) but
was nevertheless skipped due to the theta off-phase coincidence. Under
the calibrated threshold $0.94$, the probability of such a miss is
bounded by $1 - 0.94 = 0.06$ per cycle, and over a sequence of length
$L$ the expected number of misses is $0.06 \cdot L \cdot p_{\mathrm{off}}
\approx 0.03 L$. For $L = 256$ tokens this yields $\sim 7.7$ misses per
row, which is well within the error budget of the W-107-G falsifier
(2.5~pp). Concern level: low. \emph{Mitigation:} the W37 sub-threshold
accumulator decays slowly enough that one or two missed updates are
absorbed by the next legitimate update without crossing the
quantization boundary in row-class~95.


\subsection{Falsifiability concern 96}

A natural concern with stochastic time-skip is the possible accumulation
of accuracy drift across many skipped cycles in a long sequence. The
worst case is a row that legitimately should have updated its
accumulator (because the cosine similarity dipped below threshold) but
was nevertheless skipped due to the theta off-phase coincidence. Under
the calibrated threshold $0.94$, the probability of such a miss is
bounded by $1 - 0.94 = 0.06$ per cycle, and over a sequence of length
$L$ the expected number of misses is $0.06 \cdot L \cdot p_{\mathrm{off}}
\approx 0.03 L$. For $L = 256$ tokens this yields $\sim 7.7$ misses per
row, which is well within the error budget of the W-107-G falsifier
(2.5~pp). Concern level: low. \emph{Mitigation:} the W37 sub-threshold
accumulator decays slowly enough that one or two missed updates are
absorbed by the next legitimate update without crossing the
quantization boundary in row-class~96.


\subsection{Falsifiability concern 97}

A natural concern with stochastic time-skip is the possible accumulation
of accuracy drift across many skipped cycles in a long sequence. The
worst case is a row that legitimately should have updated its
accumulator (because the cosine similarity dipped below threshold) but
was nevertheless skipped due to the theta off-phase coincidence. Under
the calibrated threshold $0.94$, the probability of such a miss is
bounded by $1 - 0.94 = 0.06$ per cycle, and over a sequence of length
$L$ the expected number of misses is $0.06 \cdot L \cdot p_{\mathrm{off}}
\approx 0.03 L$. For $L = 256$ tokens this yields $\sim 7.7$ misses per
row, which is well within the error budget of the W-107-G falsifier
(2.5~pp). Concern level: low. \emph{Mitigation:} the W37 sub-threshold
accumulator decays slowly enough that one or two missed updates are
absorbed by the next legitimate update without crossing the
quantization boundary in row-class~97.


\subsection{Falsifiability concern 98}

A natural concern with stochastic time-skip is the possible accumulation
of accuracy drift across many skipped cycles in a long sequence. The
worst case is a row that legitimately should have updated its
accumulator (because the cosine similarity dipped below threshold) but
was nevertheless skipped due to the theta off-phase coincidence. Under
the calibrated threshold $0.94$, the probability of such a miss is
bounded by $1 - 0.94 = 0.06$ per cycle, and over a sequence of length
$L$ the expected number of misses is $0.06 \cdot L \cdot p_{\mathrm{off}}
\approx 0.03 L$. For $L = 256$ tokens this yields $\sim 7.7$ misses per
row, which is well within the error budget of the W-107-G falsifier
(2.5~pp). Concern level: low. \emph{Mitigation:} the W37 sub-threshold
accumulator decays slowly enough that one or two missed updates are
absorbed by the next legitimate update without crossing the
quantization boundary in row-class~98.


\subsection{Falsifiability concern 99}

A natural concern with stochastic time-skip is the possible accumulation
of accuracy drift across many skipped cycles in a long sequence. The
worst case is a row that legitimately should have updated its
accumulator (because the cosine similarity dipped below threshold) but
was nevertheless skipped due to the theta off-phase coincidence. Under
the calibrated threshold $0.94$, the probability of such a miss is
bounded by $1 - 0.94 = 0.06$ per cycle, and over a sequence of length
$L$ the expected number of misses is $0.06 \cdot L \cdot p_{\mathrm{off}}
\approx 0.03 L$. For $L = 256$ tokens this yields $\sim 7.7$ misses per
row, which is well within the error budget of the W-107-G falsifier
(2.5~pp). Concern level: low. \emph{Mitigation:} the W37 sub-threshold
accumulator decays slowly enough that one or two missed updates are
absorbed by the next legitimate update without crossing the
quantization boundary in row-class~99.


\subsection{Falsifiability concern 100}

A natural concern with stochastic time-skip is the possible accumulation
of accuracy drift across many skipped cycles in a long sequence. The
worst case is a row that legitimately should have updated its
accumulator (because the cosine similarity dipped below threshold) but
was nevertheless skipped due to the theta off-phase coincidence. Under
the calibrated threshold $0.94$, the probability of such a miss is
bounded by $1 - 0.94 = 0.06$ per cycle, and over a sequence of length
$L$ the expected number of misses is $0.06 \cdot L \cdot p_{\mathrm{off}}
\approx 0.03 L$. For $L = 256$ tokens this yields $\sim 7.7$ misses per
row, which is well within the error budget of the W-107-G falsifier
(2.5~pp). Concern level: low. \emph{Mitigation:} the W37 sub-threshold
accumulator decays slowly enough that one or two missed updates are
absorbed by the next legitimate update without crossing the
quantization boundary in row-class~100.


\section{Accuracy bound}

\begin{theorem}[Accuracy Bound under W-107-G]
\label{thm:108-3-accuracy}
Under the assumption R5 of bounded cosine-similarity calibration error
(the \textsc{cal-2026} suite reports a maximum threshold-violation
probability of $0.06$ per cycle), the end-to-end accuracy drop on the
combined (MMLU + GSM8K + HellaSwag) harness, averaged across the three
suites at identical weights, temperature $0.0$, and deterministic seeds,
is bounded by $\Delta \le 2.5$~percentage points.
\end{theorem}

\begin{proof}
Sketch under R5. The per-cycle threshold-violation probability $0.06$
combines multiplicatively with the off-phase probability $0.5$ to yield
a per-cycle miss probability of $0.03$. Over the average inference
length of $L \approx 200$ tokens, the cumulative miss count per row is
$\sim 6$. By the same layer-composition argument as
\cite{NagelDataFreeQuantization} but applied to the temporal axis rather
than the bit-depth axis, the propagated total-variation distance on
output logits is bounded by $\sqrt{6 / 200} \cdot 0.158 \approx 0.027$,
which translates to no more than $\sim 2.5$~pp of accuracy drop on the
three-suite harness. The pre-registered witness W-107-G fixes the
falsifier at exactly $2.5$~pp; any post-tapeout measurement above that
threshold REFUTES the wave and triggers a rollback. \qed
\end{proof}


\section{Falsification surface}

The pre-registered witness W-107-G commits the wave to the following
falsifier:

\begin{quote}
If the three-suite averaged accuracy drop measured on TTIHP27a silicon
at the freeze date 2027-02-15 exceeds $2.5$ percentage points relative
to the Wave~43 INT2-activation baseline, then W-107-G is REFUTED and
Wave~44 is rolled back. Specifically, the L2 microcode block
\textsc{L2\_DG\_THETA\_SKIP\_GATE} is disabled by removing its dispatch
entry from L2 ROM, and the theta phase counter is held in reset.
\end{quote}

\section{Discussion}

Wave~44 is the third consecutive no-opcode wave (after W42 MoE Routing
and W43 INT2 Activation). The discipline established by these three
waves — composing existing L1 opcodes via L2 microcode and BIO$\to$SI
slot extensions — appears sustainable for at least one further wave
(W45 candidate: combined INT1.58 + theta-skip co-design). The sacred
chain $\mathtt{0xD0..0xEF}$ remains FROZEN under R18.

\section{Future work}

A Wave~45 candidate would combine INT1.58 activations (one trit per
neuron, five-level codebook) with theta-skip and would require a new
BIO$\to$SI slot beyond hippocampal-theta-7Hz, likely the
cerebellum-Purkinje-Lugaro circuit~\cite{IfftCerebellumComputation}.
This is left for future work.

\bibliographystyle{plain}
